%%%%%%%% ACL SUBMISSION FILE %%%%%%%%%%%%%%%%%

\documentclass[11pt]{article}

% Change "review" to "final" for the camera-ready version, or to
% "preprint" for a non-anonymous version with page numbers.
\usepackage[final]{acl}

% Standard ACL package includes.
\usepackage{times}
\usepackage{latexsym}
\usepackage[T1]{fontenc}
\usepackage[utf8]{inputenc}
\DeclareUnicodeCharacter{2248}{\ensuremath{\approx}}
\DeclareUnicodeCharacter{2192}{\ensuremath{\rightarrow}}
\DeclareUnicodeCharacter{2011}{-}
\DeclareUnicodeCharacter{2013}{--}
\DeclareUnicodeCharacter{2014}{---}
\DeclareUnicodeCharacter{2018}{`}
\DeclareUnicodeCharacter{2019}{'}
\DeclareUnicodeCharacter{201C}{``}
\DeclareUnicodeCharacter{201D}{''}
\DeclareUnicodeCharacter{30FB}{\textperiodcentered}
\usepackage{microtype}
\usepackage{inconsolata}

% Paper-specific packages.
\usepackage{graphicx}
\usepackage{subcaption}
\usepackage{booktabs}
\usepackage{amsmath}
\usepackage{amssymb}
\usepackage{mathtools}
\usepackage{amsthm}
\usepackage[capitalize,noabbrev]{cleveref}
\usepackage{listings}
\usepackage{xcolor}

\lstset{
     basicstyle=\ttfamily\small,
     breaklines=true,
     frame=single,
     backgroundcolor=\color{gray!5}
}

%%%%%%%%%%%%%%%%%%%%%%%%%%%%%%%%
% THEOREMS
%%%%%%%%%%%%%%%%%%%%%%%%%%%%%%%%
\theoremstyle{plain}
\newtheorem{theorem}{Theorem}[section]
\newtheorem{proposition}[theorem]{Proposition}
\newtheorem{lemma}[theorem]{Lemma}
\newtheorem{corollary}[theorem]{Corollary}
\theoremstyle{definition}
\newtheorem{definition}[theorem]{Definition}
\newtheorem{assumption}[theorem]{Assumption}
\theoremstyle{remark}
\newtheorem{remark}[theorem]{Remark}

% Todonotes is useful during development; simply uncomment the next line
%    and comment out the line below the next line to turn off comments
\usepackage[disable,textsize=tiny]{todonotes}
%\usepackage[textsize=tiny]{todonotes}
% Preamble
\usepackage{orcidlink}
% --------导入自定义的包
\input{usepackageDefinition}

\title{WearerText: Benchmarking Wearer-Centered Scene Text Understanding in AI-Glasses Videos}

% \author{
% Yaping Zhang$^{1,2}$\thanks{\revion{Yaping Zhang, Mei Tu, and Jinting Wu contributed equally to this work.}}\,\orcidlink{0000-0001-6892-905X},
% Mei Tu$^{3}$\footnotemark[1]\,\orcidlink{0009-0005-6458-6973},
% Jinting Wu$^{3}$\footnotemark[1]\,\orcidlink{0000-0001-9642-4008},
% Qixuan Zhang$^{1,2}$\,\orcidlink{0009-0002-7726-5085},
% Xiaoyu Chen$^{1,2}$\,\orcidlink{0009-0005-6017-0700},\\
% \bfseries
% Mengchao Zhang$^{3}$\,\orcidlink{0009-0007-7167-5557},
% Yuchen Huang$^{3}$\,\orcidlink{0009-0002-2764-7501},
% Junshen Chen$^{3}$\,\orcidlink{0009-0008-1757-7531},
% Fan Zhang$^{3}$\,\orcidlink{0009-0002-4157-6981},\\
% \bfseries
% Lu Xiang$^{1,2}$\thanks{\revion{Corresponding authors: Lu Xiang and Chengqing Zong.}}\,\orcidlink{0000-0002-7702-7274},
% Junnan Zhu$^{1,2}$\,\orcidlink{0000-0002-9856-2946},
% Yang Zhao$^{1,2}$\,\orcidlink{0000-0003-1028-3406},
% Yu Zhou$^{1,2}$\,\orcidlink{0000-0002-4911-4717},
% Chengqing Zong$^{1,2}\footnotemark[2]$\orcidlink{0000-0002-9864-3818}
% \\[2mm]
% \small $^{1}$ Institute of Automation, Chinese Academy of Sciences; $^{2}$School of Artificial Intelligence,
% University of Chinese Academy of Sciences;
% \\
% \small  $^{3}$SRCB, Samsung.
% \\
% \small  \{yaping.zhang,lu.xiang, yang.zhao, yzhou,cqzong\}@nlpr.ia.ac.cn, \{mei.tu,jinting01.wu,zhang.fan,mc.zhang\}@samsung.com
% }

\author{
Yaping Zhang$^{1,2}$\,\orcidlink{0000-0001-6892-905X},
Mei Tu$^{3}$\,\orcidlink{0009-0005-6458-6973},
Jinting Wu$^{3}$\,\orcidlink{0000-0001-9642-4008},
Qixuan Zhang$^{1,2}$\,\orcidlink{0009-0002-7726-5085},
Xiaoyu Chen$^{1,2}$\,\orcidlink{0009-0005-6017-0700},\\
\bfseries
Mengchao Zhang$^{3}$\,\orcidlink{0009-0007-7167-5557},
Yuchen Huang$^{3}$\,\orcidlink{0009-0002-2764-7501},
Junshen Chen$^{3}$\,\orcidlink{0009-0008-1757-7531},
Fan Zhang$^{3}$\,\orcidlink{0009-0002-4157-6981},\\
\bfseries
Lu Xiang$^{1,2}$\,\orcidlink{0000-0002-7702-7274},
Junnan Zhu$^{1,2}$\,\orcidlink{0000-0002-9856-2946},
Yang Zhao$^{1,2}$\,\orcidlink{0000-0003-1028-3406},
Yu Zhou$^{1,2}$\,\orcidlink{0000-0002-4911-4717},
Chengqing Zong$^{1,2}$\thanks{\revion{Corresponding authors: Chengqing Zong.}}，\,\orcidlink{0000-0002-9864-3818}
\\[2mm]
\small $^{1}$ Institute of Automation, Chinese Academy of Sciences; $^{2}$School of Artificial Intelligence,
University of Chinese Academy of Sciences;
\\
\small  $^{3}$SRCB, Samsung.
\\
\small  \{yaping.zhang, lu.xiang, junnan.zhu, yang.zhao, yzhou, cqzong\}@nlpr.ia.ac.cn, \{mei.tu, jinting01.wu, zhang.fan, mc.zhang\}@samsung.com
}


\begin{document}
\maketitle

% \begin{abstract}
%      Recent MLLMs have shown potential in scene-text-aware video
% understanding, but their ability in AI-glasses scenarios remains unclear.
% AI-glasses videos are passively captured without a camera-frame preview, so the wearer cannot adjust framing as their head moves. Relevant text may therefore drift off-center or appear only briefly, while the recorded view may not reflect the wearer's attention. The challenge is therefore not merely to read text, but to identify and read the right text based on the wearer's position, viewpoint trajectory, and goal.
% To systematically evaluate this capability, we introduce
% \textbf{WearerText}, a benchmark for \textbf{Wearer}-Centered Scene
% \textbf{Text} Understanding, comprising 1{,}101 real-world AI-glasses
% videos and 15{,}336 QA pairs spanning 13 tasks. The tasks are organized
% into a three-level diagnostic hierarchy: \textit{L1 Text Perception},
% \textit{L2 Contextual Understanding}, and \textit{L3 Wearer-Grounded
% Reasoning}, targeting unstable text reading, wearer-view alignment,
% and goal-oriented text reasoning, respectively. To ensure reliable
% construction and evaluation, we propose a
% multi-agent--hybrid verification pipeline for data construction, and
% a verified-evidence metric (VerEval) for open-ended answer
% evaluation. Benchmarking 20 recent MLLMs shows that current models still
% struggle with wearer-centered scene text understanding, especially on tasks
% requiring wearer-view alignment and goal-oriented reasoning.\footnote{https://github.com/AprilYapingZhang/WearerText}   
% \end{abstract}


\begin{abstract}
    Recent MLLMs have shown promise in scene-text-aware video understanding, yet their ability to understand scene text from the wearer's perspective in AI-glasses videos remains underexplored. AI-glasses videos are passively captured without a camera-frame preview, limiting the wearer's ability to adjust framing during recording. As a result, relevant text may be unstable, and the recorded view may not fully reflect the wearer's attention. The challenge is thus not merely to read visible text, but to identify and read the right text given the wearer's position, viewpoint trajectory, and goal.
To systematically evaluate this capability, we introduce \textbf{WearerText}, a benchmark for \textbf{Wearer}-Centered Scene \textbf{Text} Understanding, comprising 1{,}101 real-world AI-glasses videos and 15{,}336 QA pairs spanning 13 tasks.
These tasks are organized into a three-level diagnostic hierarchy: \textit{L1 Text Perception}, \textit{L2 Contextual Understanding}, and \textit{L3 Wearer-Grounded Reasoning}, targeting unstable text perception, wearer-view alignment, and goal-oriented text reasoning, respectively.
To ensure reliable construction and evaluation, we develop \textsc{MAH-V}, a multi-agent hybrid verification pipeline for QA construction, and propose VerEval, a Verified-evidence metric for open-ended Evaluation.
Benchmarking 20 recent MLLMs shows that current models still struggle with wearer-centered scene text understanding, particularly on tasks requiring wearer-view alignment and goal-oriented text reasoning.\footnote{https://github.com/AprilYapingZhang/WearerText}   
\end{abstract}

%AI-glasses videos are passively captured without a camera-frame preview, so natural head motion often makes text off-center. The challenge is therefore not just reading text, but identifying the right text for the wearer's position, viewpoint trajectory, and goal. 

% \input{sec_intro_icml_v5}
\input{sec_intro_acl_v4}

%%%-----------------------分层任务定义表格------------%
\begin{table*}[t]
  \centering
  \small
  \caption{Hierarchy of Wearer-Centered Scene Text Understanding Tasks and Examples in {WearerText}.}
  \label{tab:ar-text-tasks}
  \setlength{\tabcolsep}{.4em}
  \fontsize{8.5}{10.5}\selectfont
  \resizebox{.95\linewidth}{!}{
    % \begin{tabular}{@{} c l p{0.38\linewidth} p{0.42\linewidth} @{}}
    \begin{tabular}{@{} c@{\hspace{0.38em}} l p{0.38\linewidth} p{0.38\linewidth} @{}}
      \toprule
      \textbf{Task ID} & \textbf{Task Name} & \textbf{Task Definition} & \textbf{Examples} \\
      \midrule

      \rowcolor{ModelGreen}
      \multicolumn{4}{c}{
      \textbf{Level 1: Text Perception.}
      \textit{Can the model \textsc{SEE} and \textsc{LOCALIZE} unstable visible text in dynamic AI-glasses videos?}
      } \\
      \midrule
      L1.1 & \makecell[l]{Dynamic Text\\Recognition (DTR)}
    %   L1.1   & Dynamic Text Recognition
           & Transcribe scene text from unstable egocentric views with head motion and viewpoint changes.
           & ``I was walking fast and missed the phone number on that red sign. Can you tell me what it was?'' \\
      L1.2 & \makecell[l]{Dynamic Text\\Localization (DTL)}
           & Locate queried text regions in dynamic frames and describe where they appear in the observed scene.
           & ``Where is the floor and unit number located on the billboard?'' \\

      \midrule
      \rowcolor{ModelGreen}
      \multicolumn{4}{c}{
      \textbf{Level 2: Contextual Understanding.}
      \textit{Can the model \textsc{GROUND} relevant text in the wearer's surrounding context?}
      } \\
      \midrule

      L2.1 & \makecell[l]{Query-Guided\\Text Extraction (QTE)}
           & Select and extract relevant information from dense environmental text according to the user's query.
           & ``How many years has this `Dianmao' restaurant been specializing in Peking Duck?'' \\

      L2.2 & Confirmation (CF)
           & Verify the user's prior belief or memory using current visual-textual evidence.
           & ``The sign says `Floor 4'. I'm in the mall now; this shop is on the 4th floor, right?'' \\

      L2.3 & \makecell[l]{Cross-Lingual\\Understanding (CLU)}
           & Interpret abbreviations, symbols, pinyin, or multilingual identifiers within the physical context.
           & ``What does `DIANMAO' at the top of the sign represent in this context?'' \\

      L2.4 & \makecell[l]{Situated Semantic\\Interpretation (SSI)}
           & Interpret the semantic meaning of slogans, signs, or brand descriptions beyond literal recognition.
           & ``What does the slogan `Unmatched in Beijing, Renowned Everywhere' imply about this brand?'' \\

      L2.5 & \makecell[l]{Wearer-Centered\\Situational Inference (WSI)}
           & Infer scene-level attributes or venue suitability by grounding text in the wearer's current situation.
           & ``Based on the style and the `20-year' claim, is this a good place to treat friends to authentic local food?'' \\

      \midrule
      \rowcolor{ModelGreen}
      \multicolumn{4}{c}{
      \textbf{Level 3: Wearer-Grounded Reasoning.}
      \textit{Can the model \textsc{REASON} and \textsc{ASSIST} according to viewpoint trajectories and user goals?}
      } \\
      \midrule

      L3.1 & \makecell[l]{History-Conditioned\\Intent Inference (HII)}
           & Infer the wearer's current intent by combining previous scene-text observations with the current view.
           & ``I'm tempted by the duck but have a meeting soon. Are there any `fast-eating' sets mentioned?'' \\

      L3.2 & \makecell[l]{Action-Trajectory\\Intent Prediction (AIP)}
           & Interpret the wearer's view trajectory to predict latent goals or possible next actions.
           & ``I looked at the coffee shop sign and then the arrow sign. Am I trying to find the coffee shop?'' \\

      L3.3 & \makecell[l]{Structured Text\\Synthesis (STS)}
           & Integrate structured textual sources, such as menus, or floor boards, while preserving their layout logic.
           & ``Translate the entire list of stores on this directory into English for me.'' \\

      L3.4 & \makecell[l]{Spatio-Temporal\\Navigation Reasoning (STN)}
           & Provide action guidance by integrating current location, target text, and movement cues over time.
           & ``I see the `Floor 4' ad now. What is the fastest way for me to get to the actual store from here?'' \\

      L3.5 & \makecell[l]{Goal-Oriented\\ Comparison (GOC)}
           & Compare text-described options across the video and select the one that matches the wearer's goal.
           & ``Considering both `Professionalism' and `Price', is this shop a better choice than a fast-food outlet?'' \\

      L3.6 & \makecell[l]{Viewpoint Trajectory\\Description (VTD)}
           & Describe how the wearer's viewpoint changes over time relative to text-bearing signs or directories.
           & ``Describe how my viewpoint changed.'' \\
      \bottomrule
    \end{tabular}
  }
\end{table*}







\section{Related Work}


% ~~~\textbf{Egocentric Video Understanding.} Egocentric vision has received growing attention due to its promise for real-world AI assistance. Large-scale efforts such as HD-EPIC~\cite{perrett2025hdepic}, EmbodiedBench~\cite{yang2025embodiedbench}, Ego4D~\cite{9879279} and Project Aria~\cite{engel2023project} provide long-form first‑person videos with rich annotations for activity understanding and human–object interaction. 
% These corpora are valuable for generic egocentric video understanding, but they do not specifically target scene-text understanding, wearer-view alignment, or goal-oriented text evidence reasoning, which are essential for 
% wearer-centered guidance.


% ~~~\noindent\textbf{\revion{Visual Text Understanding.}}
% Visual text understanding has advanced from robust recognition and scene-text VQA~\cite{zhang2021robust,8953586,biten2019scene} to layout-aware document analysis, image translation, and multimodal benchmarks across multilingual, scientific, and wearable settings~\cite{zhang2026promptdla,zhang2025reading,liang2025ssr,zhang2025icdardimt,li2026mmtit,zhang2025hiscibench,chang2025wearvqa}. 
% WearVQA~\cite{chang2025wearvqa}, the most related wearable benchmark, focuses on general VQA over single wearable images. 
% WearerText instead evaluates dynamic text evidence in AI-glasses videos, including cross-frame grounding, wearer-view alignment, temporal dynamics, and goal-conditioned reasoning (Table~1).


% \textbf{Text-Centric Video Understanding.}
% Text in videos carries explicit and structured semantics, and several
% benchmarks have been proposed to evaluate text-aware video
% understanding~\cite{zhao2022towards,wu2021bilingual,reddy2020roadtext,zhou2025egotextvqa,wu2023icdar}.
% RoadTextVQA~\cite{reddy2020roadtext} focuses on driving scenes with
% vehicle-mounted viewpoints, while M4-ViTeVQA~\cite{zhao2022towards}
% collects diverse YouTube videos to study general video text
% understanding. More recent benchmarks such as
% MME-VideoOCR~\cite{shi2025mmevideoocr} and
% VidText~\cite{yang2025vidtext} further evaluate OCR-related abilities
% in videos. EgoTextVQA~\cite{zhou2025egotextvqa} moves closer to
% egocentric assistance by using head-mounted videos from driving and
% indoor activities. However, these benchmarks do not systematically
% isolate the key challenges of AI-glasses text understanding caused by natural head motion, wearer-view alignment, and goal-oriented text evidence reasoning. In addition, most existing datasets rely on manual construction and standard string matching or LLM-judge evaluation, which may limit scalability and fine-grained diagnosis.

% Table~\ref{tab:comparison} compares \dataset~with existing text-aware
% and egocentric video benchmarks. Unlike prior datasets, \dataset~is
% built from real-world AI-glasses videos and is designed around a
% three-level diagnostic hierarchy that directly targets unstable text
% reading, wearer-view alignment, and goal-oriented text evidence reasoning.



\noindent\textbf{Visual Text Understanding.}
Visual text understanding has evolved from robust scene text recognition and VQA~\cite{zhang2021robust,8953586,biten2019scene} to layout-aware document analysis, image translation, and multimodal benchmarks across multilingual, scientific, and wearable settings~\cite{zhang2026promptdla,zhang2025reading,liang2025ssr,zhang2025icdardimt,li2026mmtit,zhang2025hiscibench,chang2025wearvqa}.
Among wearable benchmarks, WearVQA~\cite{chang2025wearvqa} is the most closely related, but focuses on general VQA over single wearable images.
In contrast, \dataset~targets scene text understanding in dynamic AI-glasses videos, where models must handle unstable text, align textual evidence with the wearer's view, and reason according to the wearer's goal.

\noindent\textbf{Egocentric Video Understanding.}
Egocentric benchmarks study visual understanding from a first-person perspective and are closely related to wearable assistance.
WearVQA~\cite{chang2025wearvqa} evaluates wearable-image understanding, while EgoTextVQA~\cite{zhou2025egotextvqa} extends text-aware VQA to head-mounted videos from driving and indoor activities.
However, existing egocentric benchmarks do not explicitly decompose scene text understanding into unstable text perception, wearer-view alignment, and goal-oriented reasoning over changing viewpoints.
\dataset~focuses on this wearer-centered setting, where the relevance of textual evidence depends on the wearer's position, viewpoint trajectory, and goal.

\noindent\textbf{Text-Centric Video Understanding.}
Several benchmarks evaluate text-aware video understanding~\cite{zhao2022towards,wu2021bilingual,reddy2020roadtext,zhou2025egotextvqa,wu2023icdar}.
RoadTextVQA~\cite{reddy2020roadtext} focuses on driving scenes from vehicle-mounted cameras, while M4-ViTeVQA~\cite{zhao2022towards} covers diverse YouTube videos.
More recent benchmarks, including MME-VideoOCR~\cite{shi2025mmevideoocr} and VidText~\cite{yang2025vidtext}, further evaluate OCR-related capabilities in videos.
These benchmarks primarily assess text recognition and understanding from the recorded video itself, but do not systematically evaluate whether models can identify text relevant to the wearer as the viewpoint changes.

Table~\ref{tab:comparison} summarizes these differences.
Unlike prior text-aware and egocentric benchmarks, \dataset~is built from real-world AI-glasses videos and organizes evaluation into a three-level diagnostic hierarchy.

%---------------主实验表格：数据方法介绍部分----------------------%
\begin{figure*}[t]
\centering
\includegraphics[width=0.95\textwidth]{figures/figure_methods.pdf}
\caption{
Overview of the WearerText construction pipeline.
We first define a three-level hierarchy for AI-glasses text understanding, then collect videos and generate layout-aware reference text and hybrid QA pairs.
The multi-agent--hybrid verification pipeline then performs data quality control and refinement before final human validation.
}
\label{fig:pipeline}
\end{figure*}

\section{\dataset}
\label{sec:methodology}

 
 In this section, we present the design rationale and construction pipeline for {WearerText}, which supports scalable construction and verifiable annotation for wearer-centered scene text understanding.
 

\subsection{Hierarchical Task Design}
\label{sec:task_design}

To systematically evaluate AI-glasses text understanding, we organize
\dataset~into a three-level hierarchy with \tasknum tasks
(Table~\ref{tab:ar-text-tasks}). The hierarchy follows the key
capabilities discussed above: unstable text reading, wearer-view
alignment, and goal-oriented text evidence reasoning.
\textbf{L1: Text Perception} evaluates whether models can recognize and
localize unstable or briefly visible text in dynamic AI-glasses videos.
\textbf{L2: Contextual Understanding} tests whether models can ground relevant text within the wearer's surrounding context,
including dense signage, multilingual text, slogans, and scene-level
semantics. \textbf{L3: Wearer-Grounded Reasoning} further examines
whether models can reason over viewpoint trajectories and temporal text
evidence to support goal-oriented assistance, such as navigation and viewpoint-motion understanding.
Together, these levels provide a diagnostic framework for identifying
where current MLLMs fail across perception, contextual
grounding, and wearer-centered reasoning.
% \subsection{MAH-V Data synthesis framework}

\subsection{Data Construction Pipeline}
\label{sec:data_pipeline}

To support scalable and grounded benchmark construction, 
we design an integrated pipeline that includes video collection, 
layout-aware reference text grounding, and hierarchical question--answer generation, 
as illustrated in Figure~\ref{fig:pipeline}.

\noindent\textbf{AI-glasses Video Collection}
To capture the structural complexity of real-world interactions, 
we utilize RayNeo AI Glasses to acquire a raw video dataset $\mathcal{V}$ consisting of 1,101 egocentric sequences. Each sequence, recorded at 30 FPS with a duration of 5--10 seconds, 
captures short interaction primitives such as approaching an object or scanning a menu. Our reasoning tasks therefore target goal-oriented selection of text evidence within these clips. 
For each video $v \in \mathcal{V}$, we simultaneously record a corresponding high-resolution global image $I_g$ to serve as a stable spatial reference for ground-truth text extraction. 
Detailed scenario coverage, sub-domains, and wearable-video condition statistics are provided in Appendix~\ref{sec:dataset_statistics} and Table~\ref{tab:detailed_data_stats}.

\noindent\textbf{Layout-aware reference text grounding.}
Given the global image $I_g$, we construct a reference text $T_{\text{ref}}$ to serve as the factual anchor for subsequent question generation and verification. Specifically, we first apply OCR to obtain a raw token set $T_{\text{raw}} = \text{OCR}(I_g)$. To preserve the semantic reading order, we employ an LLM-based layout-aware transformation that reorganizes $T_{\text{raw}}$ into $T_{\text{layout}}$ using spatial coordinates $(x, y, w, h)$.
Finally, human verification is performed to produce a high-quality “gold-standard” reference text $T_{\text{ref}}$.

\noindent\textbf{Hybrid hierarchical QA generation.}
Based on $(v, T_{\text{ref}})$ pairs and the predefined three-level task hierarchy, we generate question--answer instances as candidate pairs $(q_i, a_i)$, where $q_i$ denotes a wearer-centered question and $a_i$ denotes its reference answer. We use a hybrid strategy, $\mathcal{D}_{\text{hybrid}} = \mathcal{D}_{\text{auto}} \cup \mathcal{D}_{\text{man}}$. The automatic branch $\mathcal{D}_{\text{auto}}$ leverages MLLMs for scalable synthesis, while the manual branch $\mathcal{D}_{\text{man}}$ provides expert-curated instances to complement automated generation with greater diversity and quality assurance. This hybrid design aims to balance large-scale coverage with reasoning-oriented question diversity.



\subsection{MAH-V: Multi-Agent--Hybrid Verification}

The reliability of task construction is critical for evaluating situated reasoning. We introduce the \textbf{M}ulti-\textbf{A}gent--\textbf{H}ybrid \textbf{V}erification (MAH-V) pipeline, which combines a committee of specialized agents $\mathcal{A} = \{A_{\text{fact}}, A_{\text{logic}}, A_{\text{consistency}}, A_{\text{refine}}\}$ with a final human-verification stage to filter $\mathcal{D}_{\text{hybrid}}$. The agent committee runs the multi-agent quality control stage shown in Figure~\ref{fig:pipeline}, and its output is reviewed by expert annotators as described at the end of this section.

% \noindent\textbf{Fact Checker Agent.}
% The Fact Checker $A_{\text{fact}}$ validates incorrect information by ensuring strict accuracy. For a candidate pair $(q_i, a_i) \in \mathcal{D}_{\text{hybrid}}$, $A_{\text{fact}}$ takes the video $v$ and reference text $T_{\text{ref}}$ together as inputs to compute a factuality score $S_{fact} \in [0, 1]$:
% \begin{equation}
% S_{fact} = A_{\text{fact}}(q_i, a_i, v, T_{\text{ref}})
% \end{equation}
% It verifies whether numerical values or entities in $a_i$ exist within the verified $T_{\text{ref}}$. Instances where $S_{fact} < \tau$ (threshold) are marked for improvement.


\noindent\textbf{Fact Checker Agent.}
\revion{The Fact Checker $A_{\text{fact}}$ assesses the factual accuracy of each candidate pair $(q_i,a_i)\in\mathcal{D}_{\text{hybrid}}$. Given the video $v$ and the verified reference text $T_{\text{ref}}$, it produces a structured factual assessment:}
\begin{equation}
\revion{S_{fact}=A_{\text{fact}}(q_i,a_i,v,T_{\text{ref}}).}
\end{equation}
\revion{Here, $S_{fact}=(s_{fact},\mathcal{R}_{fact})$, where $s_{fact}\in[0,1]$ is the factuality score and $\mathcal{R}_{fact}$ contains concrete revision suggestions. The agent checks whether factual elements in the question and answer, particularly numerical values and entities, are supported by the video and $T_{\text{ref}}$. Candidates with $s_{fact}<\tau$ are marked for refinement, while $\mathcal{R}_{fact}$ identifies unsupported content and provides actionable suggestions for correcting the question or answer.}



\noindent\textbf{Logic Checker Agent.}
The Logic Checker $A_{\text{logic}}$ evaluates whether the reasoning process is grounded in the wearer-centered video context along three criteria.
First, it assesses sequence necessity, determining whether temporal visual evidence is required to answer the question.
Second, it performs indexical alignment by linking spatial references such as ``the sign on the left'' to specific video segments, ensuring that spatial descriptions are grounded in the wearer's viewpoint.
Finally, it verifies temporal causality by checking that the textual cues supporting the answer have appeared in the observed video evidence, preventing reasoning that relies on unavailable information.

\noindent\textbf{Consistency Agent.}
To improve the correctness of the generated data, the {Consistency Agent} $A_{\text{consistency}}$ performs cross validation. It detects contradictions between the factual evidence retrieved by $A_{\text{fact}}$ and the reasoning logic proposed by $A_{\text{logic}}$. For example, if $A_{\text{fact}}$ confirms a bus departs at ``10:30'' but $A_{\text{logic}}$ infers a wait time based on ``10:35'', $A_{\text{consistency}}$ identifies the mismatch and generates a conflict report $\mathcal{R}_{\text{conf}}$.

\noindent\textbf{Refinement Agent.}
The {Refinement Agent} $A_{\text{refine}}$ serves as the final improvement step. Rather than simply discarding failed pairs, it uses the descriptive feedback and conflict reports $\{\mathcal{R}_{\text{conf}}, S_{fact}\}$ to improve the final QA pairs: $(q_i', a_i') = A_{\text{refine}}(q_i, a_i, \mathcal{R}_{\text{conf}}, S_{fact})$.


\noindent\textbf{Human Verification and Final Dataset.}
After multi-agent verification, the refined QA pairs are passed to expert human annotators for final validation. Three domain experts independently review each sample, assessing both factual accuracy and logical coherence. Only samples achieving unanimous approval or majority consensus (2 out of 3) are retained in the final benchmark. This human-in-the-loop validation achieves an inter-annotator agreement of $\kappa = 0.89$, indicating high consistency in our quality assessment protocol.


\subsection{VerEval: Verified-Evidence Evaluation}
We propose \textbf{Ver}ified-Evidence \textbf{Eval}uation (\textbf{VerEval}), which adapts the MAH-V agent architecture for open-ended answer evaluation. While MAH-V uses OCR reference text for quality control during benchmark construction, VerEval directly compares model predictions with ground truth.
VerEval uses three scoring agents. The \textit{Fact Scorer} produces \(S_{fact}\) by checking whether factual anchors such as numbers, timestamps, and entities match the visual evidence and ground truth. The \textit{Logic Scorer} produces \(S_{logic}\) by evaluating whether the answer's reasoning is compatible with spatio-temporal constraints in the video. The \textit{Consistency Scorer} then produces \(S_{consist}\) by jointly reviewing the fact and logic judgments, their justifications, and the model answer. Thus, \(S_{consist}\) is the output of the consistency-scoring prompt, not a renamed factuality score.

We use \(S_{fact}\) as an intermediate diagnostic signal, but exclude it from the final aggregate because its evidence is already checked by the consistency scorer and because the held-out human-alignment study showed lower discrimination for the raw fact score than for \(S_{consist}\) and \(S_{logic}\). The final VerEval score is therefore
\begin{equation}
S_{overall} = \lambda_1 \cdot S_{consist} + \lambda_2 \cdot S_{logic}.
\end{equation}
In this paper, we set \(\lambda_1=0.9\) and \(\lambda_2=0.1\), selected on held-out human-judged samples; detailed weighting results are provided in Appendix~\ref{sec:vereval_weighting_ablation}. We report the alignment between VerEval and human judgment in Figure~\ref{fig:score_distribution}, and provide all three scoring prompts in Appendix~\ref{scoring_prompt}.


%---------------主实验表格：Main experimental results----------------------%
% \begin{table*}[t]
% \caption{Overall Performance Results on \dataset Three-Level Task Hierarchy. The best and second-best ANLS /SAS results are \textbf{bolded} and \underline{underlined} respectively}
% \label{tab:overall_results}
% \setlength{\tabcolsep}{.2em}
% \fontsize{9}{11}\selectfont
% \resizebox{\textwidth}{!}{
% \begin{threeparttable}
% \begin{tabular}{l|cc|cccc|ccccc}
% \Xhline{1pt}
% \textbf{Model} 
%   & \multicolumn{2}{c|}{\textbf{L1: Core Cognition}} 
%   & \multicolumn{4}{c|}{\textbf{L2: Contextual Understanding}} 
%   & \multicolumn{5}{c}{\textbf{L3: Wearer-Grounded Reasoning}} \\
% \Xhline{1pt}
% \rowcolor{ModelGreen}
%  & \textbf{DTR} & \textbf{TG} 
%    & \textbf{DQ} & \textbf{CF} & \textbf{CLU} & \textbf{GTT} 
%    & \textbf{MID} & \textbf{GT}  & \textbf{SI}  & \textbf{STR} & \textbf{CS} \\
% \Xhline{1pt}
% \multicolumn{10}{l}{\emph{\textcolor{lightgray}{Closed-source Models}}} \\
% GPT-4o & 26.9/33 & 8.3/14 & 19/25.5 & 8.5/14.7 & 19.6/27.6 & 33.6/37.3 & \underline{28.1}/35.9 & 2.4/4.8 & 12.6/19.6 & 58.5/63.6 & 16.2/23.6 \\
% o4-mini-high & 33.6/39.6 & 9.4/16.7 & \textbf{33.4}/40.9 & 17.5/27.5 & 23.1/31.6 & \textbf{42}/\textbf{43.1} & 27.1/34.0 & 2.7/4.4 & \textbf{18.6}/\textbf{29.1} & \underline{59}/\underline{63.9} & \textbf{41.3}/\textbf{49.9} \\
% Gemini-2.5-Flash & 28.1/34.5 & 9.4/14.1 & 25.7/33.4 & 9.9/17.2 & 17.2/23.8 & 32.9/37 & 22.1/29.9 & 1/2 & 12.1/18 & 51.5/55.9 & \underline{35.7}/\underline{43.4} \\
% Gemini-2.5-Pro & \underline{35.8}/\underline{41.9} & 9.8/15.5 & 25.6/33.5 & 15.5/23.3 & 19.9/26.5 & \underline{39.7}/\underline{40.5} & 18.1/25.1 & 4.9/9.1 & 11.3/18.8 & 54.4/59.1 & 20.8/32.6 \\
% Qwen-VL-Max & \textbf{38.4}/\textbf{44.7} & 1.9/7.9 & 28.3/37.1 & 15.6/23.6 & 11.7/15.9 & 35.9/38.4 & 11.8/16.2 & 1.6/2.5 & \underline{17.9}/\underline{25.5} & \textbf{60.7}/\textbf{65.1} & 17.4/24.3 \\
% Qwen-VL-OCR & 6.2/10.4 & 1.8/2.7 & 1.6/2.9 & 1.1/1.9 & 6.3/6.9 & 9.9/10.4 & - & 0.1/0.1 & 1.4/2 & 1.3/2.2 & 3.1/5.1 \\
% \hline
% \multicolumn{10}{l}{\emph{\textcolor{lightgray}{Open-source Models}}} \\
% InternVL3-8B & 13.9/21.5 & 6.1/12.2 & 27.6/37.2 & 16.8/24.3 & 24.6/34.4 & 20.4/21.8 & 25.6/35.8 & 15.3/23.5 & 10.3/20.1 & 21.4/29.7 & 20.7/33.7 \\
% InternVL3-38B & 18.1/23.1 & 48.6/55.3 & \underline{32.2}/\textbf{43.5} & \textbf{28}/\textbf{39.3} & 18.2/26.7 & 27.1/30 & 20.4/29.1 & \underline{39.9}/\underline{46} & 11.6/19 & 44/49.5 & 23.8/34.6 \\
% InternVL3-78B & 19.8/24.9 & \textbf{49.2}/\textbf{55.8} & 21.2/30.3 & 14.1/22.4 & 26.8/36.8 & 28.1/29.6 & 28.0/\textbf{37.5} & 39.6/45.4 & 9.5/18.4 & 45.2/50.6 & 22.2/36.7 \\
% Qwen2.5-VL-3B & 17.8/22.8 & 37/43.7 & 17.9/25.7 & 18.1/26.7 & \underline{28}/\underline{37.3} & 24.5/27.6 & 15.2/20.6 & 26.4/37.6 & 8.8/13.9 & 37.6/41.8 & 13.4/20.1 \\
% Qwen2.5-VL-7B & 13.8/18.5 & 36.4/43.5 & 29.7/40.8 & 15.1/22.4 & 20.9/29.6 & 25/29.3 & \textbf{28.4}/\underline{37.1} & 35.6/44.3 & 13.2/21.2 & 37/41.8 & 23.5/31.9 \\
% Qwen2.5-VL-32B & 21/26.5 & 45.7/52.2 & 11.5/16.6 & 6/9.5 & 23/32.6 & 29.8/31.5 & 22.1/29.0 & 34.5/41.3 & 3.9/6.1 & 42.1/47.2 & 5.5/8.4 \\
% Qwen2.5-VL-72B & 22.2/27.5 & \underline{48.9}/\underline{55.5} & 13.5/19.5 & 8.2/13.3 & \textbf{29.9}/\textbf{40.3} & 29.4/31 & 19.7/26.0 & \textbf{43}/\textbf{50.7} & 5.7/9 & 46.4/51.4 & 8.5/14 \\
% \hline
% \multicolumn{10}{l}{\emph{\textcolor{lightgray}{Open-source Video Models}}} \\
% VideoLLaMA3-7B & 1.5/9.3 & 0.0/0.0 & 2.6/21.8 & 9.6/\underline{28.8} & 3.2/22.6 & 1.0/15.3 & 4.6/30.4 & 0.0/8.2 & 1.4/17.8 & 4.1/14.0 & 0.2/21.6 \\
% LLaVA-NeXT-Video-7B & 12.1/23.8 & 0.0/0.0 & 9.9/25.7 & 5.9/24.6 & 5.5/27.2 & 9.0/30.8 & 1.3/13.9 & 0.0/12.2 & 5.2/21.8 & 45.7/56.6 & 2.5/18.9 \\
% Valley & 18.2/24.1 & 10.6/14.9 & 28.8/\underline{41.0} & \underline{19.6}/28.5 & 26.7/35.2 & 5.2/6.6 & 20.0/27.4 & 7.3/11.1 & 14.9/25.1 & 41.0/47.1 & 30.0/39.4 \\
% \Xhline{1pt}
% \end{tabular}
% \end{threeparttable}
%  }
% \vspace{-0.4cm}
% \end{table*}




% \begin{table*}[t]
% \caption{Overall Performance Results on \dataset  Three-Level Task Hierarchy. The best and second-best logical\_avg /consistency\_avg results are \textbf{bolded} and \underline{underlined} respectively}
% \label{tab:overall_results}
% \setlength{\tabcolsep}{.2em}
% \fontsize{9}{11}\selectfont
% \resizebox{\textwidth}{!}{
% \begin{threeparttable}
% \begin{tabular}{l|cc|ccccc|cccccc}
% \Xhline{1pt}
% \textbf{Model} & \multicolumn{2}{c|}{\textbf{L1: Core Cognition}} & \multicolumn{5}{c|}{\textbf{L2: Contextual Understanding}} & \multicolumn{6}{c}{\textbf{L3: Wearer-Grounded Reasoning}} \\
% \Xhline{1pt}
% \rowcolor{ModelGreen}
%  & \textbf{DTR} & \textbf{TG} & \textbf{DQ} & \textbf{CF} & \textbf{CLU} & \textbf{GTT} & \textbf{NEW} & \textbf{MID} & \textbf{GT} & \textbf{SI} & \textbf{STR} & \textbf{CS} & \textbf{NEW} \\
% \Xhline{1pt}
% \multicolumn{14}{l}{\emph{\textcolor{lightgray}{Closed-source Models}}} \\
% GPT-5.2 & -- & -- & -- & -- & -- & -- & -- & -- & -- & -- & -- & -- & -- \\
% Gemini-3-Pro & -- & -- & -- & -- & -- & -- & -- & -- & -- & -- & -- & -- & -- \\
% Claude-Opus-4.5 & -- & -- & -- & -- & -- & -- & -- & -- & -- & -- & -- & -- & -- \\
% Qwen-VL-Max & -- & -- & -- & -- & -- & -- & -- & -- & -- & -- & -- & -- & -- \\
% \hline
% \multicolumn{14}{l}{\emph{\textcolor{lightgray}{Open-source Models}}} \\
% InternVL3-8B & -- & -- & -- & -- & -- & -- & -- & -- & -- & -- & -- & -- & -- \\
% InternVL3-38B & -- & -- & -- & -- & -- & -- & -- & -- & -- & -- & -- & -- & -- \\
% InternVL3-78B & -- & -- & -- & -- & -- & -- & -- & -- & -- & -- & -- & -- & -- \\
% Qwen2.5-VL-3B & 60.7/61.5 & 46.9/48.6 & 47.9/48.8 & 51.6/54.4 & 64.1/66.2 & 32.1/32.3 & 16.9/15.0 & 53.1/55.3 & 54.3/58.8 & 62.1/62.8 & 46.3/48.5 & 48.5/48.6 & 9.0/13.4 \\
% Qwen2.5-VL-7B & 71.7/73.6 & 57.0/62.8 & 56.5/56.0 & 62.9/64.8 & 71.6/71.7 & 39.2/39.5 & 14.9/14.8 & 53.6/56.7 & 56.9/63.8 & 60.1/64.2 & 46.5/51.2 & 46.5/49.8 & \textbf{34.7}/\textbf{41.8} \\
% Qwen2.5-VL-32B & 71.2/73.5 & 61.9/64.9 & 56.6/54.0 & \underline{68.4}/\underline{67.0} & 73.6/73.9 & 40.6/42.2 & \underline{33.2}/\textbf{29.6} & \underline{62.5}/\underline{63.2} & \underline{71.6}/\underline{70.7} & \underline{70.9}/\underline{72.0} & \underline{60.7}/\underline{58.6} & \underline{61.9}/\underline{63.0} & 15.5/22.0 \\
% Qwen3-VL-4B & 71.3/72.8 & 55.2/57.9 & 54.4/53.1 & 62.8/60.0 & 73.0/71.8 & 37.7/38.8 & 22.0/20.5 & 59.1/58.1 & 56.0/58.7 & 62.4/64.4 & 49.5/50.6 & 54.6/48.2 & 15.9/20.2 \\
% Qwen3-VL-8B & \underline{75.1}/\underline{74.6} & \underline{64.3}/\underline{66.8} & 57.7/54.9 & 57.5/61.2 & \underline{74.7}/\underline{75.3} & 44.7/43.5 & 29.0/25.8 & 47.9/50.3 & 53.2/57.2 & 51.8/57.9 & 40.6/42.4 & 39.1/42.3 & 23.1/18.9 \\
% Qwen3-VL-32B & \textbf{75.9}/\textbf{76.5} & \textbf{69.5}/\textbf{70.2} & \textbf{67.5}/\textbf{61.6} & \textbf{72.4}/\textbf{67.8} & \textbf{82.7}/\textbf{77.5} & \textbf{60.2}/\textbf{55.6} & \textbf{33.6}/\underline{27.7} & \textbf{78.4}/\textbf{74.0} & \textbf{75.5}/\textbf{71.9} & \textbf{80.6}/\textbf{77.0} & \textbf{65.4}/\textbf{61.8} & \textbf{70.8}/\textbf{68.1} & \underline{24.0}/\underline{28.2} \\
% GLM-4.6V-Flash & 55.6/58.6 & 50.3/57.7 & 56.5/58.2 & 44.0/52.8 & 50.9/51.3 & 32.9/36.5 & 20.1/19.1 & 30.8/36.5 & 27.8/35.9 & 34.9/42.0 & 30.4/37.1 & 29.2/35.6 & 8.4/11.4 \\
% GLM-4.6V & 71.1/73.9 & 60.5/65.1 & \underline{61.1}/\underline{61.3} & 56.0/58.8 & 65.8/67.9 & \underline{47.0}/\underline{47.6} & 29.4/27.0 & 43.1/48.9 & 41.8/53.1 & 46.8/56.1 & 39.0/47.7 & 35.6/40.1 & 20.6/25.8 \\
% \hline
% \multicolumn{14}{l}{\emph{\textcolor{lightgray}{Open-source Video Models}}} \\
% VideoLLaMA3-7B & -- & -- & -- & -- & -- & -- & -- & -- & -- & -- & -- & -- & -- \\
% LLaVA-NeXT-Video-7B & -- & -- & -- & -- & -- & -- & -- & -- & -- & -- & -- & -- & -- \\
% VideoR1-7B & -- & -- & -- & -- & -- & -- & -- & -- & -- & -- & -- & -- & -- \\
% Molmo-7B & -- & -- & -- & -- & -- & -- & -- & -- & -- & -- & -- & -- & -- \\
% \Xhline{1pt}
% \end{tabular}
% \end{threeparttable}
%  }
% \vspace{-0.4cm}
% \end{table*}


% \begin{table*}[t]
% \caption{Overall Performance. Scores are calculated as $0.1 \times \text{Logic} + 0.9 \times \text{Consistency}$. \textbf{Bold} indicates the best, \underline{Underline} indicates the second best.}
% \label{tab:final_overall_results}
% \centering
% \resizebox{\textwidth}{!}{
% \begin{tabular}{l | cc | ccccc | cccccc | c}\Xhline{1pt}
% \textbf{Model} & \multicolumn{2}{c|}{\textbf{L1: Core}} & \multicolumn{5}{c|}{\textbf{L2: Contextual}} & \multicolumn{6}{c|}{\textbf{L3: Wearer-Grounded Reasoning}} & \textbf{Overall} \\
%  & \textbf{DTR} & \textbf{TG} & \textbf{DQ} & \textbf{CF} & \textbf{CLU} & \textbf{GTT} & \textbf{NEW} & \textbf{MID} & \textbf{GT} & \textbf{SI} & \textbf{STR} & \textbf{CS} & \textbf{NEW} & \textbf{AVG} \\
% \Xhline{1pt}
% \multicolumn{15}{l}{\emph{Closed-source Models}} \\
% GPT-5.2 & 84.5 & \textbf{77.7} & 72.4 & 74.2 & \textbf{85.8} & \textbf{73.0} & 45.1 & \textbf{81.2} & \textbf{82.6} & \textbf{82.4} & \textbf{73.6} & \textbf{83.2} & 34.4 & \cellcolor{gray!15}{\textbf{73.1}} \\
% Gemini-3-Pro & \textbf{91.4} & 75.8 & \textbf{82.7} & \textbf{82.8} & \underline{83.2} & \underline{70.2} & \textbf{46.2} & \underline{77.1} & \underline{79.5} & 81.6 & \underline{73.1} & \underline{80.6} & 23.0 & \cellcolor{gray!15}{\underline{72.9}} \\
% Claude-Opus-4.5 & 86.9 & 71.6 & \underline{76.2} & \underline{77.4} & 72.3 & 60.5 & 44.8 & 75.2 & 79.2 & 80.1 & 71.3 & 75.7 & 35.9 & \cellcolor{gray!15}{69.8} \\
% Qwen3-VL-PLUS & \underline{89.1} & \underline{76.7} & 73.3 & 74.0 & 79.6 & 66.4 & 43.8 & 73.9 & 78.1 & \underline{82.4} & 70.4 & 78.6 & \underline{38.9} & \cellcolor{gray!15}{71.2} \\
% Seed-1.6 & 84.0 & 70.4 & 71.8 & 75.0 & 61.1 & 63.3 & \underline{46.0} & 58.6 & 54.3 & 60.2 & 55.4 & 57.6 & 35.1 & \cellcolor{gray!15}{61.0} \\
% \hline
% \multicolumn{15}{l}{\emph{Open-source Models}} \\
% InternVL3-8B & 54.4 & 49.6 & 43.9 & 50.4 & 57.0 & 28.1 & 7.9 & 40.4 & 25.5 & 49.6 & 43.3 & 35.3 & 21.4 & \cellcolor{gray!15}{39.0} \\
% InternVL3-38B & 68.5 & 54.8 & 38.3 & 53.6 & 54.2 & 32.8 & 16.2 & 42.3 & 36.3 & 50.5 & 42.6 & 44.3 & 11.2 & \cellcolor{gray!15}{42.0} \\
% Qwen2.5-VL-3B & 65.9 & 52.1 & 48.2 & 59.5 & 61.8 & 29.3 & 15.2 & 49.0 & 47.9 & 51.6 & 41.0 & 44.1 & 7.5 & \cellcolor{gray!15}{44.1} \\
% Qwen2.5-VL-7B & 73.4 & 62.2 & 56.0 & 64.6 & 71.7 & 39.4 & 14.8 & 56.4 & 63.1 & 63.8 & 50.8 & 49.5 & \textbf{41.1} & \cellcolor{gray!15}{54.4} \\
% Qwen2.5-VL-32B & 73.2 & 64.6 & 54.3 & 67.2 & 73.9 & 42.1 & 30.0 & 63.1 & 70.8 & 71.9 & 58.8 & 62.9 & 21.4 & \cellcolor{gray!15}{58.0} \\
% Qwen3-VL-4B & 72.6 & 57.7 & 53.3 & 60.2 & 72.0 & 38.7 & 20.6 & 58.2 & 58.4 & 64.2 & 50.5 & 48.8 & 19.8 & \cellcolor{gray!15}{51.9} \\
% Qwen3-VL-8B & 74.6 & 66.6 & 55.2 & 60.8 & 75.2 & 43.6 & 26.2 & 50.1 & 56.8 & 57.3 & 42.3 & 42.0 & 19.3 & \cellcolor{gray!15}{51.5} \\
% Qwen3-VL-32B & 76.5 & 70.1 & 62.2 & 68.3 & 78.1 & 56.1 & 28.3 & 74.5 & 72.2 & 77.4 & 62.1 & 68.3 & 27.8 & \cellcolor{gray!15}{63.2} \\
% GLM-4.6V-Flash-9B & 58.3 & 56.9 & 58.0 & 51.9 & 51.3 & 36.1 & 19.2 & 36.0 & 35.0 & 41.3 & 36.5 & 35.0 & 11.1 & \cellcolor{gray!15}{40.5} \\
% GLM-4.6V-106B-A12B & 73.7 & 64.6 & 61.2 & 58.5 & 67.7 & 47.6 & 27.2 & 48.3 & 52.0 & 55.1 & 46.8 & 39.6 & 25.2 & \cellcolor{gray!15}{51.4} \\
% MINICPM-V-4.5 & 76.5 & 63.3 & 62.7 & 69.6 & 46.6 & 48.4 & 22.2 & 53.5 & 50.6 & 60.5 & 46.2 & 46.0 & 9.5 & \cellcolor{gray!15}{50.4} \\
% \hline
% \multicolumn{15}{l}{\emph{Open-source Video Models}} \\
% VideoLLaMA3-7B & 51.1 & 55.0 & 47.1 & 52.8 & 53.8 & 28.5 & 13.9 & 43.1 & 50.0 & 49.0 & 39.7 & 42.7 & 4.9 & \cellcolor{gray!15}{40.9} \\
% LLaVA-NeXT-Video & 8.6 & 18.3 & 10.6 & 18.3 & 32.9 & 7.5 & 2.6 & 17.0 & 19.8 & 21.6 & 15.1 & 21.5 & 2.1 & \cellcolor{gray!15}{15.1} \\
% VideoR1-7B & 66.0 & 51.0 & 42.7 & 58.4 & 66.4 & 38.3 & 24.2 & 55.7 & 62.4 & 61.6 & 49.6 & 54.4 & 24.6 & \cellcolor{gray!15}{50.4} \\
% MOLMO-2-8B & 27.0 & 51.0 & 30.8 & 48.6 & 63.0 & 22.1 & 11.3 & 34.5 & 46.9 & 43.7 & 35.7 & 35.8 & 11.9 & \cellcolor{gray!15}{35.6} \\
% \hline
% \Xhline{1pt}
% \end{tabular}
% }
% \end{table*}



% \begin{table*}[t]
% \caption{Overall Performance. Scores are calculated as $0.1 \times \text{Logic} + 0.9 \times \text{Consistency}$. \textbf{Bold} indicates the best, \underline{Underline} indicates the second best.}
% \label{tab:final_overall_results}
% \centering
% \resizebox{\textwidth}{!}{
% \begin{tabular}{l | cc | ccccc | cccccc | c}\Xhline{1pt}
% % \textbf{Model} & \multicolumn{2}{c|}{\textbf{L1: Perception}} & \multicolumn{5}{c|}{\textbf{L2: Understanding}} & \multicolumn{6}{c|}{\textbf{L3: Wearer-Grounded Reasoning}} & \textbf{Overall} \\
% % \Xhline{1pt}

% %  & \textbf{L1.1} & \textbf{L1.2} & \textbf{L2.1} & \textbf{L2.2} & \textbf{L2.3} & \textbf{L2.4} & \textbf{L2.5} & \textbf{L3.1} & \textbf{L3.2} & \textbf{L3.3} & \textbf{L3.4} & \textbf{L3.5} & \textbf{L3.6} & \textbf{AVG} \\

% \multicolumn{1}{c|}{\textbf{Model}} & \multicolumn{2}{c|}{\textbf{L1: Perception}} & \multicolumn{5}{c|}{\textbf{L2: Understanding}} & \multicolumn{6}{c|}{\textbf{L3: Wearer-Grounded Reasoning}} & \textbf{Overall} \\

%  & \textbf{L1.1} & \textbf{L1.2} & \textbf{L2.1} & \textbf{L2.2} & \textbf{L2.3} & \textbf{L2.4} & \textbf{L2.5} & \textbf{L3.1} & \textbf{L3.2} & \textbf{L3.3} & \textbf{L3.4} & \textbf{L3.5} & \textbf{L3.6} & \multicolumn{1}{c}{\textbf{AVG}} \\
%  & \textbf{DTR} & \textbf{TG} & \textbf{DQ} & \textbf{CF} & \textbf{CLU} & \textbf{SSI} & \textbf{SI} & \textbf{DGII} & \textbf{BBIP} & \textbf{SBT} & \textbf{STR} & \textbf{CS} & \textbf{VTD} & \\

% \Xhline{1pt}
% \multicolumn{15}{l}{\emph{Closed-source Models}} \\
% GPT-5.2 & 84.5 & \textbf{77.7} & 72.4 & 74.2 & \textbf{85.8} & \textbf{73.0} & \textbf{82.4} & \textbf{81.2} & \textbf{82.6} & 45.1 & \textbf{73.6} & \textbf{83.2} & 34.4 & \cellcolor{gray!15}{\textbf{73.1}} \\
% Gemini-3-Pro & \textbf{91.4} & 75.8 & \textbf{82.7} & \textbf{82.8} & \underline{83.2} & \underline{70.2} & 81.6 & \underline{77.1} & \underline{79.5} & \textbf{46.2} & \underline{73.1} & \underline{80.6} & 23.0 & \cellcolor{gray!15}{\underline{72.9}} \\
% Claude-Opus-4.5 & 86.9 & 71.6 & \underline{76.2} & \underline{77.4} & 72.3 & 60.5 & 80.1 & 75.2 & 79.2 & 44.8 & 71.3 & 75.7 & \underline{35.9} & \cellcolor{gray!15}{69.8} \\
% Qwen3-VL-PLUS & \underline{89.1} & \underline{76.7} & 73.3 & 74.0 & 79.6 & 66.4 & \underline{82.4} & 73.9 & 78.1 & 43.8 & 70.4 & 78.6 & \textbf{38.9} & \cellcolor{gray!15}{71.2} \\
% Seed-1.6 & 84.0 & 70.4 & 71.8 & 75.0 & 61.1 & 63.3 & 60.2 & 58.6 & 54.3 & \underline{46.0} & 55.4 & 57.6 & 35.1 & \cellcolor{gray!15}{61.0} \\
% \hline
% \multicolumn{15}{l}{\emph{Open-source Models}} \\
% InternVL3-8B & 54.4 & 49.6 & 43.9 & 50.4 & 57.0 & 28.1 & 49.6 & 40.4 & 25.5 & 7.9 & 43.3 & 35.3 & 21.4 & \cellcolor{gray!15}{39.0} \\
% InternVL3-38B & 68.5 & 54.8 & 38.3 & 53.6 & 54.2 & 32.8 & 50.5 & 42.3 & 36.3 & 16.2 & 42.6 & 44.3 & 11.2 & \cellcolor{gray!15}{42.0} \\
% Qwen2.5-VL-3B & 65.9 & 52.1 & 48.2 & 59.5 & 61.8 & 29.3 & 51.6 & 49.0 & 47.9 & 15.2 & 41.0 & 44.1 & 7.5 & \cellcolor{gray!15}{44.1} \\
% % Qwen2.5-VL-7B & 73.4 & 62.2 & 56.0 & 64.6 & 71.7 & 39.4 & 63.8 & 56.4 & 63.1 & 14.8 & 50.8 & 49.5 & \textbf{41.1} & \cellcolor{gray!15}{54.4} \\
% {Qwen2.5-VL-7B} & 69.0 & 55.8 & 58.4 & 59.5 & 73.3 & 43.3 & 69.1 & 56.6 & 66.1 & 52.3 & 50.8 & 17.8 & 22.9 & \cellcolor{gray!15}{53.5}\\
% Qwen2.5-VL-32B & 73.2 & 64.6 & 54.3 & 67.2 & 73.9 & 42.1 & 71.9 & 63.1 & 70.8 & 30.0 & 58.8 & 62.9 & 21.4 & \cellcolor{gray!15}{58.0} \\
% Qwen3-VL-4B & 72.6 & 57.7 & 53.3 & 60.2 & 72.0 & 38.7 & 64.2 & 58.2 & 58.4 & 20.6 & 50.5 & 48.8 & 19.8 & \cellcolor{gray!15}{51.9} \\
% Qwen3-VL-8B & 74.6 & 66.6 & 55.2 & 60.8 & 75.2 & 43.6 & 57.3 & 50.1 & 56.8 & 26.2 & 42.3 & 42.0 & 19.3 & \cellcolor{gray!15}{51.5} \\
% Qwen3-VL-32B & 76.5 & 70.1 & 62.2 & 68.3 & 78.1 & 56.1 & 77.4 & 74.5 & 72.2 & 28.3 & 62.1 & 68.3 & 27.8 & \cellcolor{gray!15}{63.2} \\
% GLM-4.6V-Flash-9B & 58.3 & 56.9 & 58.0 & 51.9 & 51.3 & 36.1 & 41.3 & 36.0 & 35.0 & 19.2 & 36.5 & 35.0 & 11.1 & \cellcolor{gray!15}{40.5} \\
% GLM-4.6V-106B-A12B & 73.7 & 64.6 & 61.2 & 58.5 & 67.7 & 47.6 & 55.1 & 48.3 & 52.0 & 27.2 & 46.8 & 39.6 & 25.2 & \cellcolor{gray!15}{51.4} \\
% MINICPM-V-4.5 & 76.5 & 63.3 & 62.7 & 69.6 & 46.6 & 48.4 & 60.5 & 53.5 & 50.6 & 22.2 & 46.2 & 46.0 & 9.5 & \cellcolor{gray!15}{50.4} \\
% \hline
% \multicolumn{15}{l}{\emph{Open-source Video Models}} \\
% VideoLLaMA3-7B & 51.1 & 55.0 & 47.1 & 52.8 & 53.8 & 28.5 & 49.0 & 43.1 & 50.0 & 13.9 & 39.7 & 42.7 & 4.9 & \cellcolor{gray!15}{40.9} \\
% LLaVA-NeXT-Video & 8.6 & 18.3 & 10.6 & 18.3 & 32.9 & 7.5 & 21.6 & 17.0 & 19.8 & 2.6 & 15.1 & 21.5 & 2.1 & \cellcolor{gray!15}{15.1} \\
% VideoR1-7B & 66.0 & 51.0 & 42.7 & 58.4 & 66.4 & 38.3 & 61.6 & 55.7 & 62.4 & 24.2 & 49.6 & 54.4 & 24.6 & \cellcolor{gray!15}{50.4} \\
% MOLMO-2-8B & 27.0 & 51.0 & 30.8 & 48.6 & 63.0 & 22.1 & 43.7 & 34.5 & 46.9 & 11.3 & 35.7 & 35.8 & 11.9 & \cellcolor{gray!15}{35.6} \\
% \hline
% \Xhline{1pt}
% \end{tabular}
% }
% \end{table*}

\begin{table*}[t]
\caption{Overall VerEval scores. \textbf{Bold} indicates the best, \underline{Underline} indicates the second best.}
\label{tab:final_overall_results}
\centering
\resizebox{\textwidth}{!}{
\begin{tabular}{l | cc | ccccc | cccccc | c}\Xhline{1pt}
\textbf{Model} & \multicolumn{2}{c|}{\textbf{L1: Perception}} & \multicolumn{5}{c|}{\textbf{L2: Understanding}} & \multicolumn{6}{c|}{\textbf{L3: Wearer-Grounded Reasoning}} & \textbf{Overall} \\
 & \textbf{L1.1} & \textbf{L1.2} & \textbf{L2.1} & \textbf{L2.2} & \textbf{L2.3} & \textbf{L2.4} & \textbf{L2.5} & \textbf{L3.1} & \textbf{L3.2} & \textbf{L3.3} & \textbf{L3.4} & \textbf{L3.5} & \textbf{L3.6} & \textbf{AVG} \\
 & \textbf{DTR} & \textbf{DTL} & \textbf{QTE} & \textbf{CF} & \textbf{CLU} & \textbf{SSI} & \textbf{WSI} & \textbf{HII} & \textbf{AIP} & \textbf{STS} & \textbf{STN} & \textbf{GOC} & \textbf{VTD} & \\
\Xhline{1pt}
\multicolumn{15}{l}{\emph{Closed-source Models}} \\
GPT-5.2 & 84.5 & \textbf{77.7} & 72.4 & 74.2 & \textbf{85.8} & \textbf{73.0} & \textbf{82.4} & \textbf{81.2} & \textbf{82.6} & 45.1 & \textbf{73.6} & \textbf{83.2} & 34.4 & \cellcolor{gray!15}{\textbf{73.1}} \\
Gemini-3-Pro & \textbf{91.4} & 75.8 & \textbf{82.7} & \textbf{82.8} & \underline{83.2} & \underline{70.2} & 81.6 & \underline{77.1} & \underline{79.5} & \textbf{46.2} & \underline{73.1} & \underline{80.6} & 23.0 & \cellcolor{gray!15}{\underline{72.9}} \\
Claude-Opus-4.5 & 86.9 & 71.6 & \underline{76.2} & \underline{77.4} & 72.3 & 60.5 & 80.1 & 75.2 & 79.2 & 44.8 & 71.3 & 75.7 & \underline{35.9} & \cellcolor{gray!15}{69.8} \\
Qwen3-VL-PLUS & \underline{89.1} & \underline{76.7} & 73.3 & 74.0 & 79.6 & 66.4 & \underline{82.4} & 73.9 & 78.1 & 43.8 & 70.4 & 78.6 & \textbf{38.9} & \cellcolor{gray!15}{71.2} \\
Seed-1.6 & 84.0 & 70.4 & 71.8 & 75.0 & 61.1 & 63.3 & 60.2 & 58.6 & 54.3 & \underline{46.0} & 55.4 & 57.6 & 35.1 & \cellcolor{gray!15}{61.0} \\
\hline
\multicolumn{15}{l}{\emph{Open-source Models}} \\
InternVL3.5-8B & 54.4 & 49.6 & 43.9 & 50.4 & 57.0 & 28.1 & 49.6 & 40.4 & 25.5 & 7.9 & 43.3 & 35.3 & 21.4 & \cellcolor{gray!15}{39.0} \\
InternVL3.5-38B & 68.5 & 54.8 & 38.3 & 53.6 & 54.2 & 32.8 & 50.5 & 42.3 & 36.3 & 16.2 & 42.6 & 44.3 & 11.2 & \cellcolor{gray!15}{42.0} \\
Qwen2.5-VL-3B & 65.9 & 52.1 & 48.2 & 59.5 & 61.8 & 29.3 & 51.6 & 49.0 & 47.9 & 15.2 & 41.0 & 44.1 & 7.5 & \cellcolor{gray!15}{44.1} \\
{Qwen2.5-VL-7B} & 69.0 & 55.8 & 58.4 & 59.5 & 73.3 & 43.3 & 69.1 & 56.6 & 66.1 & 52.3 & 50.8 & 17.8 & 22.9 & \cellcolor{gray!15}{53.5}\\
Qwen2.5-VL-32B & 73.2 & 64.6 & 54.3 & 67.2 & 73.9 & 42.1 & 71.9 & 63.1 & 70.8 & 30.0 & 58.8 & 62.9 & 21.4 & \cellcolor{gray!15}{58.0} \\
Qwen3-VL-4B & 72.6 & 57.7 & 53.3 & 60.2 & 72.0 & 38.7 & 64.2 & 58.2 & 58.4 & 20.6 & 50.5 & 48.8 & 19.8 & \cellcolor{gray!15}{51.9} \\
Qwen3-VL-8B & 74.6 & 66.6 & 55.2 & 60.8 & 75.2 & 43.6 & 57.3 & 50.1 & 56.8 & 26.2 & 42.3 & 42.0 & 19.3 & \cellcolor{gray!15}{51.5} \\
Qwen3-VL-32B & 76.5 & 70.1 & 62.2 & 68.3 & 78.1 & 56.1 & 77.4 & 74.5 & 72.2 & 28.3 & 62.1 & 68.3 & 27.8 & \cellcolor{gray!15}{63.2} \\
GLM-4.6V-Flash-9B & 58.3 & 56.9 & 58.0 & 51.9 & 51.3 & 36.1 & 41.3 & 36.0 & 35.0 & 19.2 & 36.5 & 35.0 & 11.1 & \cellcolor{gray!15}{40.5} \\
GLM-4.6V-106B-A12B & 73.7 & 64.6 & 61.2 & 58.5 & 67.7 & 47.6 & 55.1 & 48.3 & 52.0 & 27.2 & 46.8 & 39.6 & 25.2 & \cellcolor{gray!15}{51.4} \\
MINICPM-V-4.5 & 76.5 & 63.3 & 62.7 & 69.6 & 46.6 & 48.4 & 60.5 & 53.5 & 50.6 & 22.2 & 46.2 & 46.0 & 9.5 & \cellcolor{gray!15}{50.4} \\
\hline
\multicolumn{15}{l}{\emph{Open-source Video Models}} \\
VideoLLaMA3-7B & 51.1 & 55.0 & 47.1 & 52.8 & 53.8 & 28.5 & 49.0 & 43.1 & 50.0 & 13.9 & 39.7 & 42.7 & 4.9 & \cellcolor{gray!15}{40.9} \\
LLaVA-NeXT-Video & 8.6 & 18.3 & 10.6 & 18.3 & 32.9 & 7.5 & 21.6 & 17.0 & 19.8 & 2.6 & 15.1 & 21.5 & 2.1 & \cellcolor{gray!15}{15.1} \\
VideoR1-7B & 66.0 & 51.0 & 42.7 & 58.4 & 66.4 & 38.3 & 61.6 & 55.7 & 62.4 & 24.2 & 49.6 & 54.4 & 24.6 & \cellcolor{gray!15}{50.4} \\
MOLMO-2-8B & 27.0 & 51.0 & 30.8 & 48.6 & 63.0 & 22.1 & 43.7 & 34.5 & 46.9 & 11.3 & 35.7 & 35.8 & 11.9 & \cellcolor{gray!15}{35.6} \\
\hline
\Xhline{1pt}
\end{tabular}
}
\end{table*}





\section{Experiments}
% This section details the experimental evaluation of state-of-the-art models on \dataset. We define the evaluation metrics, describe the baseline models selected, and detail the experimental setup.

\noindent\textbf{Data Statistics.}
We summarize dataset characteristics---scenario coverage, QA and video language distribution, and the test/train split---in Appendix Table~\ref{tab:detailed_data_stats}, which illustrates the dataset's diversity across scenes and text settings.

\noindent\textbf{Model Evaluation.}
%-----ToDO proof valley 和GLM 引用
We evaluate 5 closed-source API-based models (GPT-5.2~\cite{gpt-5.2}, Gemini-3-Pro~\cite{gemini-3-pro}, Claude-Opus-4.5~\cite{claude-opus-4.5}, Qwen3-VL-PLUS~\cite{Qwen3-VL}, Seed-1.6~\cite{doubao1.6}) and 15 advanced open-source MLLMs (Intern-VL3.5-8B/38B~\cite{wang2025internvl3_5}, Qwen2.5-VL-3B/7B/32B~\cite{bai2025qwen2}, {Qwen3-VL-4B/8B/32B}~\cite{Qwen3-VL}, {GLM-4.6V-Flash-9B, GLM-4.6V-106B-A12B}\cite{vteam2025glm45vglm41vthinkingversatilemultimodal}, {VideoLLaMA3-7B} \cite{zhang2025videollama}, MINICPM-V-4.5~\cite{yu2025minicpmv45cookingefficient}, LLaVA-NeXT-Video-7B~\cite{zhang2024llavanextvideo}, VideoR1-7B~\cite{feng2025video}, MOLMO-2-8B~\cite{clark2026molmo2}).
% -----TODO 确认这几个模型在最新benchmark里面的效果----%。
%The models have reported the \todo{state-of-the-art (SOTA)} on existing scene-text VQA benchmarks \cite{biten2019scene,8953586,8578478}, \eg, Qwen2-VL achieves QA accuracy $\sim$85\% on ST-VQA \cite{biten2019scene}. 
%~\ref{sup:models} and~\ref{sup:prompt}.

\subsection{Main Results}




\noindent\textbf{Alignment Between VerEval and Human Judgment.}
We evaluate VerEval on 390 QA pairs annotated by three experts and finalized by majority vote as \textit{Correct} or \textit{Wrong}. We compare against SAS, a transformer-based semantic answer similarity metric~\citep{risch2021semantic}, and ANLS, a normalized edit-similarity metric widely used in scene-text VQA~\citep{biten2019scene}. As shown in Figure~\ref{fig:score_distribution}, our VerEval Score aligns more closely with human judgments. It yields a larger Correct--Wrong mean gap (0.402 vs. 0.112/0.094). These results indicate that VerEval provides a more reliable evaluation signal than ANLS and SAS for open-ended wearer-centered scene text understanding. We therefore report the main leaderboard with VerEval scores in Table~\ref{tab:final_overall_results}, while detailed ANLS and SAS results are provided in Appendix Tables~\ref{tab:anls_results} and~\ref{tab:sas_results}. VerEval produces scores on the raw [0,1] scale; we multiply them by 100 when reporting results for readability in Tables.



\begin{figure}[htbp]
\centering
\includegraphics[width=\columnwidth]{figures/Fig4_humanAgentAlignmentboxplots.pdf}
\caption{Distribution of automated metric scores grouped by human judgment (Wrong vs. Correct). (a) VerEval. (b) SAS. (c) ANLS. Red diamonds indicate mean values. The larger separation for VerEval indicates stronger alignment with human.}
\label{fig:score_distribution}
\end{figure}

\textls[-30]{\noindent\textbf{Capability Boundaries of Mainstream MLLMs.} As detailed in Table~\ref{tab:final_overall_results}, strong closed-source models such as GPT-5.2 and Gemini-3-Pro obtain the highest overall scores in our evaluation (73.1 and 72.9 respectively), but clear limitations remain in \textit{Structured Text Synthesis (L3.3 STS)} and \textit{Viewpoint Trajectory Description (L3.6 VTD)}. Specifically, the best scores in these two categories are 46.2 (Gemini-3-Pro) and 38.9 (Qwen3-VL-PLUS), indicating a substantial drop relative to \textit{Dynamic Text Recognition (L1.1 DTR)}. The results suggest that moving from unstable text perception to goal-oriented text evidence selection remains challenging for current MLLMs.}
\begin{figure}[t]
  \centering
  \includegraphics[width=0.78\linewidth]{figures/Fig-radar_compact.pdf}
  \caption{Capability Boundaries of Mainstream MLLMs.}
  \label{fig:radar}
  \vspace{-.2cm}
\end{figure}



\textbf{Effectiveness of Training with Scalable Verified Data.} To test whether MAH-V data is useful beyond benchmark construction, we fine-tune Qwen2.5-VL-7B on (v1) unverified LLM-generated data and (v2) MAH-V-verified data---both drawn from videos disjoint from the test split---and compare against the zero-shot baseline (v0). Table~\ref{tab:qwen_ablation} shows that v1 gains on L1/L3 but loses $1.2\%$ on L2 \textit{Contextual Understanding}, whereas v2 improves every level ($+3.7\%$ L1, $+1.8\%$ L2, $+0.5\%$ L3). Multi-agent verification therefore both raises overall data quality and removes the L2 degradation introduced by direct LLM synthesis.


\begin{figure}[t]
  \centering
  \includegraphics[width=0.95\linewidth]{figures/Fig6_ModelScaling.pdf}
\caption{Performance across model sizes and task levels. Across both Qwen2.5-VL and Qwen3-VL families, L1 remains easier than L2, and L3 is consistently the most challenging, indicating that the three-level hierarchy preserves a clear difficulty separation across scales.}
  \label{fig:model_size_scaling}
  \vspace{-.2cm}
\end{figure}


\begin{figure*}[t]
    \centering
    \includegraphics[width=0.95\linewidth]{figures/Fig7_failure_case_0526.pdf}
    \caption{
Representative failure cases and Error distribution on WearerText. Red text denotes model predictions, while green text denotes reference answers.
}
\label{fig:faifurecases}
\end{figure*}


\begin{figure}[htb]
  \centering
  \setlength{\abovecaptionskip}{0.6cm}    
\setlength{\belowcaptionskip}{-0.3cm}
  % \includegraphics[width=\linewidth]{figures/diagnostic_bars_combined.pdf}
  \includegraphics[width=\linewidth]{figures/diagnostic_bars_combined.pdf}
  \caption{Diagnostic bar comparisons for GPT-5.2. (a) Modality ablation: full video, sampled keyframes, layout-aware OCR text. (b) Memory strategies: no memory, OCR text memory, structured text memory with reasoning, and raw multimodal memory. Arrows indicate changes relative to the full-video baseline.}
  \label{fig:diagnostic_bars}
  \vspace{-.2cm}
\end{figure}





\begin{table}[t]
\caption{Effectiveness of Scalable Verified Data Generation. \textbf{Bold} indicates the best performance; \underline{underline} indicates the second best. \textcolor{blue}{$\uparrow$} and \textcolor{red}{$\downarrow$} denote improvement or decline relative to the v0 baseline.}
\label{tab:qwen_ablation}
\centering
\resizebox{.95\linewidth}{!}{
\begin{tabular}{l | c | c c c}
\Xhline{1pt}
\textbf{Metric} & \makecell{\textbf{v0}\\\textbf{Zero-shot}} & \makecell{\textbf{v1}\\\textbf{Unverified LLM}} & \makecell{\textbf{v2}\\\textbf{MAH-V}} \\
\Xhline{1pt}
\textbf{L1} & 62.4 & \underline{65.0} \textcolor{blue}{$\uparrow$2.6} & \textbf{66.1} \textcolor{blue}{$\uparrow$3.7} \\
\textbf{L2} & 60.7 & 59.5 \textcolor{red}{$\downarrow$1.2} & \textbf{62.5} \textcolor{blue}{$\uparrow$1.8} \\
\textbf{L3} & 56.5 & \textbf{59.0} \textcolor{blue}{$\uparrow$2.5} & \underline{57.0} \textcolor{blue}{$\uparrow$0.5} \\
\hline
\textbf{Overall} & 59.5 & \underline{60.3} \textcolor{blue}{$\uparrow$0.9} & \textbf{61.4} \textcolor{blue}{$\uparrow$1.9} \\
\Xhline{1pt}
\end{tabular}
}
\end{table}
\vspace{-3mm}





\subsection{Key Findings}
We conduct extensive experiments to thoroughly analyze the challenges presented by \dataset.  Our analysis centers on the following questions: 
\noindent\textbf{Q1: How do tasks differ in difficulty?}
As shown in Figure \ref{fig:model_size_scaling}, across both Qwen2.5-VL and Qwen3-VL families, task difficulty increases from L1 to L3. L1 perception scores highest, L2 understanding trails behind, and L3 reasoning remains most challenging. This clear separation ($L1 > L2 > L3$) confirms that task complexity grows with each level. Notably, the substantial performance gap between L1 and L3 (up to 20 points) suggests that current MLLMs struggle to move from reading unstable text to wearer-view alignment and goal-oriented text evidence selection.




\noindent\textbf{Q2: Why do current models fail?}
To move beyond aggregate scores, we manually categorize GPT-5.2's incorrect predictions on the held-out test set. As shown in Figure~\ref{fig:faifurecases}(d), the dominant failure type is {goal-oriented text evidence reasoning} (57.4\%), followed by {unstable text recognition} (26.6\%) and {wearer-view alignment} (16.0\%). These results suggest that the key bottleneck is not only OCR, but the ability to select and reason over the right textual evidence under the wearer's view. Figure~\ref{fig:faifurecases} illustrates these failure modes. In (a), the model fails to recognize briefly visible scrolling text on a digital screen, showing the difficulty of unstable text reading under blur and short exposure. In (b), the model incorrectly includes years from visually similar awards, indicating weak temporal evidence accumulation and fine-grained text selection. In (c), the model gives an incorrect navigation suggestion because it fails to align visible shop signs with the wearer's spatial context. Overall, these cases show that AI-glasses text understanding requires the joint ability to read unstable text, ground it from the wearer's viewpoint, and accumulate goal-relevant evidence over time. 

\noindent\textbf{Q3: Does the benchmark require full AI-glasses video?}
We compare GPT-5.2 with full video, sampled keyframes, and layout-aware OCR text. As shown in Figure~\ref{fig:diagnostic_bars}(a), full video consistently outperforms reduced inputs across L1--L3 and overall. The gap is modest for image-only input but much larger for OCR-only input, especially at L3, indicating that higher-level tasks require visual context and temporally ordered wearer-view signals beyond extracted text.



\noindent\textbf{Q4: Does memory help AI-glasses text understanding?}
We compare four memory strategies on GPT-5.2 in Figure~\ref{fig:diagnostic_bars}(b): no memory, unstructured multimodal memory with raw past frames, OCR-derived text memory with concatenated OCR tokens, and structured text memory with distilled per-frame text indexed by the wearer's viewpoint trajectory and paired with a reasoning prompt. Three task-specific patterns emerge.
(i)~\textbf{Memory helps when evidence is temporal.} Gains concentrate on L3 \textit{Wearer-Grounded Reasoning} (e.g., L3.1 HII and L3.5 GOC), while frame-local L1 perception is largely unaffected. This suggests that the L3 deficit in Q1 partly reflects a memory-architecture limitation rather than model capacity alone.
(ii)~\textbf{Structure beats raw recall.} OCR-only memory brings little gain, whereas trajectory-indexed distilled text with a reasoning prompt improves overall performance by $+2.0$. It is also the only configuration that helps L3.4 STN and L3.6 VTD, both of which explicitly query the wearer's movement.
(iii)~\textbf{Raw multimodal memory hurts.} Stale frames degrade accuracy below the no-memory baseline, suggesting that re-perception competes with current-frame perception and amplifies the unstable-text errors.
\emph{Design implication:} future systems should maintain a low-bandwidth episodic record of distilled text indexed by the viewpoint trajectory rather than a verbatim multimodal buffer.











%
\section{Conclusion}

% We introduced \dataset, a benchmark for evaluating AI-glasses text understanding in real-world wearable scenarios. Through our three-level task hierarchy, we assessed {20} recent MLLMs on real-world public text-use challenges. Our findings show that current models remain limited on wearer-centered scene text understanding and reasoning, especially when answers require selecting relevant text evidence from partial observations according to the wearer's goal. These results suggest that progress in this setting will require stronger wearer-view alignment, structured memory over observed text, and intent-aware reasoning.
We introduce \dataset, a benchmark for evaluating AI-glasses text understanding in real-world wearable scenarios. Through our three-level task hierarchy, we assess {20} recent MLLMs on real-world public text-use challenges. Our findings have shown that current models remain limited on wearer-centered scene text understanding and reasoning, especially when answers require selecting relevant text evidence from partial observations according to the wearer's goal. These results suggest that progress in this setting requires stronger wearer-view alignment, structured memory over observed text, and intent-aware reasoning.

\section*{Acknowledgments}
\revion{This work was supported by the National Natural Science Foundation of China (NSFC) under Grant No.~62476275.}

\section*{Limitations}

\dataset~ is collected from real wearable-glasses videos, but its coverage is still bounded by the recording locations, participant behaviors, device characteristics, and public-space scenarios included in our study. Although the benchmark contains multilingual and mixed-language text, its distribution is not intended to represent all languages, regions, signage conventions, or accessibility needs. Future extensions should include broader geographic coverage, more diverse users, and additional low-resource languages.

Our task design focuses on offline question answering over short egocentric clips. Thus, ``goal-oriented text evidence selection'' refers to selecting relevant textual evidence within 5--10 second interaction primitives according to the wearer's current goal, rather than long-horizon memory over minutes or hours. This setting enables controlled comparison across models, but it does not fully capture interactive use cases in which a wearable assistant can ask clarification questions, adapt its guidance, or operate under real-time latency and energy constraints.

\dataset~ is wearer-centered in the sense that questions are grounded in the camera-view trajectory of the person wearing the device. However, our RayNeo recordings do not include eye-tracking signals, so camera view should not be interpreted as true gaze. Future versions with synchronized gaze sensing could separate head direction, eye fixation, and user intention more precisely.



\section*{Ethical Considerations}

This work advances AI-glasses text understanding by providing a comprehensive and rigorous benchmark. 
 All data contributors provided informed consent. To ensure privacy, we implemented rigorous anonymization (facial blurring) for all third-party individuals. Our pipeline prioritizes transparency and bias mitigation throughout the collection and annotation process.

%\textbf{Societal Impact}. This work provides a rigorous testbed for advancing text-literate egocentric multimodal systems. It contributes to key application scenarios such as assistive navigation and embodied agents. Potential domain and language biases require attention when transferring to other application domains.


\bibliography{example_paper}

%%%%%%%%%%%%%%%%%%%%%%%%%%%%%%%%%%%%%%%%%%%%%%%%%%%%%%%%%%%%%%%%%%%%%%%%%%%%%%%
%%%%%%%%%%%%%%%%%%%%%%%%%%%%%%%%%%%%%%%%%%%%%%%%%%%%%%%%%%%%%%%%%%%%%%%%%%%%%%%
% APPENDIX
%%%%%%%%%%%%%%%%%%%%%%%%%%%%%%%%%%%%%%%%%%%%%%%%%%%%%%%%%%%%%%%%%%%%%%%%%%%%%%%
%%%%%%%%%%%%%%%%%%%%%%%%%%%%%%%%%%%%%%%%%%%%%%%%%%%%%%%%%%%%%%%%%%%%%%%%%%%%%%%
\clearpage
\appendix

%%%%%%%%%%%%%%%%%%%%%%%%%%%%%%%%%%%%%%%%%%%%%%%%%%%%%%%%%%%%%%%%%%%%%%%%%%%%%%%

\section{Dataset Statistics}
\label{sec:dataset_statistics}

We summarize the statistics of \textbf{WearerText} in Table~\ref{tab:detailed_data_stats}. The dataset consists of 1,101 real-world egocentric videos collected in three text-rich scenarios---shopping centers, transportation hubs, and tourist destinations---and contains 15,336 question--answer pairs spanning diverse perceptual, contextual, and reasoning-oriented tasks. \revion{Figure~\ref{fig:per_task_examples} presents representative qualitative examples across all 13 WearerText tasks. The examples illustrate the progression from dynamic text perception (L1), through contextual understanding (L2), to wearer-grounded reasoning over goals, history, and viewpoint trajectories (L3).}

The collection spans sub-domains such as restaurants, supermarkets, menus, subway and railway stations, airports, street scenes, parks, and museums. Common wearable-video conditions related to the challenges discussed in the main paper are intentionally preserved: 72.0\% of videos contain viewpoint changes, 33.7\% contain briefly visible text, and 4.8\% contain visible motion blur. These conditions are not mutually exclusive, since a single video may contain multiple challenges. We view this coverage as diagnostic rather than exhaustive; broader geographic and linguistic expansion remains future work.

The benchmark primarily covers Chinese and English language settings. Most QA pairs are in Chinese (67.59\%), while English and mixed Chinese-English QA pairs account for 16.40\% and 15.19\%, respectively. At the video level, 29.22\% of the visible scene text is Chinese-only, while 69.15\% contains mixed Chinese-English text, reflecting the multilingual nature of everyday embodied text environments.

% We reserve 1{,}391 QA pairs as the held-out test split used for all zero-shot model evaluations and human-alignment analysis in the main paper. The remaining 13{,}945 QA pairs form a separate train split that is reserved for the MAH-V SFT ablation (Table ~\ref{tab:qwen_ablation}) and is \emph{not} used to tune the zero-shot baselines on the leaderboard; the fine-tuned Qwen2.5-VL-7B numbers therefore should not be compared directly with the zero-shot rows in Table~\ref{tab:final_overall_results}.

\revion{We reserve 1{,}391 QA pairs as the held-out test split used for the complete 13-task zero-shot leaderboard and human-alignment analysis in the main paper. The remaining 13{,}945 QA pairs form a separate training pool and are never used to tune the zero-shot leaderboard models. The controlled Qwen2.5-VL-7B SFT comparison in Table~\ref{tab:qwen_ablation} evaluates v0, v1, and v2 on the same 11-task subset (L1.1--L3.4), excluding L3.5 and L3.6; its absolute scores are therefore not directly comparable to the complete 13-task results in Table~\ref{tab:final_overall_results}.}


\begin{table}[t]
\centering
\caption{Statistics of \textbf{WearerText}. Wearable-video condition percentages are non-exclusive: one video may contain multiple conditions.}
\label{tab:detailed_data_stats}
\small
\begin{tabular}{@{}l@{\hspace{0.8em}}r@{}}
\toprule
\textbf{Statistic} & \textbf{Value} \\
\midrule
Total Videos & 1,101 \\
Total Questions & 15,336 \\
\midrule
\multicolumn{2}{@{}l}{\textbf{By Scenario}} \\
\quad Shopping Centers & 378 videos \\
\quad Transportation Hubs & 412 videos \\
\quad Tourist Destinations & 311 videos \\
\midrule
\multicolumn{2}{@{}l}{\textbf{Wearable-Video Conditions (non-exclusive)}} \\
\quad Viewpoint Changes & 72.0\% videos \\
\quad Briefly Visible Text & 33.7\% videos \\
\quad Visible Motion Blur & 4.8\% videos \\
\midrule
\multicolumn{2}{@{}l}{\textbf{QA Language}} \\
\quad Chinese & 67.59\% \\
\quad English & 16.40\% \\
\quad Mixed Chinese-English & 15.19\% \\
\quad Other & 0.82\% \\
\midrule
\multicolumn{2}{@{}l}{\textbf{Video Text Language}} \\
\quad Chinese & 29.22\% \\
\quad Mixed Chinese-English & 69.15\% \\
\midrule
\multicolumn{2}{@{}l}{\textbf{Verification Split}} \\
\quad Test Set & 1,391 QAs / 107 videos \\
\quad Train Set & 13,945 QAs / 994 videos \\
\bottomrule
\end{tabular}
\end{table}




\begin{figure*}[t]
  \centering
  \includegraphics[width=\textwidth]{figures/Figure-Examples.pdf}
  \caption{\revion{Per-task qualitative examples from WearerText. The examples cover all 13 tasks in the three-level hierarchy and pair representative frames from AI-glasses videos with the corresponding bilingual wearer-centered queries.}}
  \label{fig:per_task_examples}
\end{figure*}






%%%%%%%%%%%%%%%%%%%%%%%%%%%%%%%%%%%%%%%%%%%%%%%%%%%%%%%%%%%%%%%%%%%%%%%%%%%%%%%
\section{Multi-Faceted Per-Task Evaluation}
\label{sec:multi_eval}

To complement VerEval, we also report two surface-similarity metrics so that readers can compare \dataset~models under traditional VQA scoring conventions:
(1) \textbf{ANLS} (Average Normalized Levenshtein Similarity) measures exact string accuracy via normalized Levenshtein similarity (edit distance divided by the longer string length), shown in Table~\ref{tab:anls_results}.
(2) \textbf{SAS} (Semantic Answer Similarity) measures the semantic similarity between the predicted answer and the ground truth answer, so paraphrased but equivalent answers still score highly, shown in Table~\ref{tab:sas_results}.

% ========================================
% ANLS Score Table
% ========================================
\begin{table*}[t]
\caption{ANLS Performance. \textbf{Bold} indicates the best, \underline{Underline} indicates the second best. Task abbreviations: DTR (Dynamic Text Recognition), DTL (Dynamic Text Localization), QTE (Query-Guided Text Extraction), CF (Confirmation), CLU (Cross-Lingual Understanding), SSI (Situated Semantic Interpretation), WSI (Wearer-Centered Situational Inference), HII (History-Conditioned Intent Inference), AIP (Action-Trajectory Intent Prediction), STS (Structured Text Synthesis), STN (Spatio-Temporal Navigation Reasoning), GOC (Goal-Oriented Comparison), VTD (Viewpoint Trajectory Description).}
\label{tab:anls_results}
\centering
\resizebox{\textwidth}{!}{
\begin{tabular}{l | cc | ccccc | cccccc | c}\Xhline{1pt}
\textbf{Model} & \multicolumn{2}{c|}{\textbf{L1: Perception}} & \multicolumn{5}{c|}{\textbf{L2: Understanding}} & \multicolumn{6}{c|}{\textbf{L3: Wearer-Grounded Reasoning}} & \textbf{Overall} \\
 & \textbf{L1.1} & \textbf{L1.2} & \textbf{L2.1} & \textbf{L2.2} & \textbf{L2.3} & \textbf{L2.4} & \textbf{L2.5} & \textbf{L3.1} & \textbf{L3.2} & \textbf{L3.3} & \textbf{L3.4} & \textbf{L3.5} & \textbf{L3.6} & \textbf{AVG} \\
 & \textbf{DTR} & \textbf{DTL} & \textbf{QTE} & \textbf{CF} & \textbf{CLU} & \textbf{SSI} & \textbf{WSI} & \textbf{HII} & \textbf{AIP} & \textbf{STS} & \textbf{STN} & \textbf{GOC} & \textbf{VTD} & \\
\Xhline{1pt}
\multicolumn{15}{l}{\emph{Closed-source Models}} \\
GPT-5.2 & \textbf{62.4} & \textbf{19.7} & \underline{20.3} & 8.6 & \textbf{16.4} & \textbf{23.0} & 15.9 & 0.5 & 0.0 & 0.6 & 0.0 & 0.0 & 0.0 & \cellcolor{gray!15}{12.9} \\
Gemini-3-Pro & 59.0 & 18.7 & \textbf{29.9} & \textbf{21.0} & 19.1 & \underline{30.3} & \textbf{20.0} & \textbf{3.2} & \textbf{0.0} & \textbf{1.5} & \textbf{0.0} & \textbf{1.6} & \underline{6.6} & \cellcolor{gray!15}{\textbf{16.2}} \\
Claude-Opus-4.5 & 68.8 & \underline{20.6} & 27.9 & \underline{9.8} & 6.5 & 18.3 & \underline{19.6} & 0.0 & 0.0 & 0.0 & 0.0 & 0.5 & 0.0 & \cellcolor{gray!15}{13.2} \\
Qwen3-VL-PLUS & \underline{73.1} & 20.0 & \underline{29.9} & 13.9 & \underline{10.5} & 21.4 & 18.2 & 0.5 & 0.0 & 0.5 & 0.0 & 0.0 & 0.0 & \cellcolor{gray!15}{\underline{14.5}} \\
Seed-1.6 & 72.4 & 17.5 & 32.9 & 10.7 & 7.5 & 19.7 & 20.7 & \underline{1.9} & 0.0 & 0.0 & 0.0 & 0.0 & \textbf{5.1} & \cellcolor{gray!15}{14.5} \\
\hline
\multicolumn{15}{l}{\emph{Open-source Models}} \\
InternVL3.5-8B & 34.6 & 2.7 & 6.6 & 3.4 & 3.7 & 3.9 & 1.1 & 0.0 & 0.0 & 0.0 & 0.0 & 0.0 & 1.2 & \cellcolor{gray!15}{4.4} \\
InternVL3.5-38B & 51.2 & 16.9 & 10.5 & 1.6 & 2.7 & 7.7 & 7.2 & 0.0 & 0.0 & 0.0 & 0.0 & 0.0 & 14.9 & \cellcolor{gray!15}{8.7} \\
Qwen2.5-VL-3B & 43.9 & 7.0 & 7.8 & 5.4 & 3.4 & 6.1 & 6.4 & 0.7 & 0.0 & 0.5 & 0.0 & 0.0 & 22.9 & \cellcolor{gray!15}{8.0} \\
Qwen2.5-VL-7B & 50.4 & 3.8 & 13.3 & 5.9 & 9.3 & 10.8 & 5.0 & 0.6 & 0.0 & 0.0 & 0.0 & 0.0 & 20.0 & \cellcolor{gray!15}{9.1} \\
Qwen2.5-VL-32B & 61.1 & 12.2 & 15.0 & 4.6 & 8.4 & 9.6 & 11.6 & 0.5 & 0.0 & 0.0 & 0.0 & 0.0 & 12.7 & \cellcolor{gray!15}{10.4} \\
Qwen3-VL-4B & 55.3 & 20.7 & 13.0 & 4.0 & 9.8 & 11.0 & 8.5 & 0.6 & 0.0 & 0.5 & 0.0 & 0.0 & 9.9 & \cellcolor{gray!15}{10.3} \\
Qwen3-VL-8B & 61.1 & 16.4 & 20.2 & 2.8 & 8.9 & 14.1 & 11.2 & 0.0 & 0.0 & 0.0 & 0.0 & 0.0 & 0.0 & \cellcolor{gray!15}{10.4} \\
Qwen3-VL-32B & 59.7 & 15.1 & 16.0 & 4.4 & 9.2 & 11.6 & 12.7 & 0.0 & 0.0 & 0.0 & 0.0 & 0.0 & 2.4 & \cellcolor{gray!15}{10.1} \\
GLM-4.6V-Flash-9B & 51.4 & 15.7 & 19.9 & 1.6 & 7.8 & 10.7 & 8.9 & 0.0 & 0.0 & 0.0 & 0.0 & 0.0 & 0.0 & \cellcolor{gray!15}{8.9} \\
GLM-4.6V-106B-A12B & 62.1 & 22.1 & 21.8 & 8.3 & 9.7 & 17.8 & 10.7 & 0.5 & 0.0 & 0.0 & 0.0 & 0.0 & 12.1 & \cellcolor{gray!15}{12.7} \\
MINICPM-V-4.5 & 51.3 & 14.7 & 23.1 & 11.2 & 3.8 & 9.6 & 6.8 & 0.0 & 0.0 & 0.0 & 0.0 & 0.0 & 0.0 & \cellcolor{gray!15}{9.3} \\
\hline
\multicolumn{15}{l}{\emph{Open-source Video Models}} \\
VideoLLaMA3-7B & 16.7 & 5.2 & 1.9 & 2.2 & 1.8 & 2.7 & 2.6 & 0.0 & 0.0 & 0.0 & 0.0 & 0.0 & 0.0 & \cellcolor{gray!15}{2.5} \\
LLaVA-NeXT-Video & 0.0 & 0.5 & 0.5 & 0.0 & 0.0 & 0.5 & 0.0 & 0.0 & 0.0 & 0.0 & 0.0 & 0.0 & 0.0 & \cellcolor{gray!15}{0.1} \\
VideoR1-7B & 1.5 & 1.2 & 2.4 & 0.0 & 1.1 & 1.7 & 0.5 & 0.0 & 0.0 & 0.0 & 0.0 & 0.0 & 0.8 & \cellcolor{gray!15}{0.7} \\
MOLMO-2-8B & 15.7 & 14.0 & 8.3 & 0.8 & 4.8 & 2.4 & 3.7 & 0.0 & 0.0 & 0.0 & 0.0 & 0.0 & 0.0 & \cellcolor{gray!15}{3.8} \\
\hline
\Xhline{1pt}
\end{tabular}
}
\end{table*}

% ========================================
% SAS Score Table
% ========================================
\begin{table*}[t]
\caption{SAS Performance. \textbf{Bold} indicates the best, \underline{Underline} indicates the second best. Task abbreviations: DTR (Dynamic Text Recognition), DTL (Dynamic Text Localization), QTE (Query-Guided Text Extraction), CF (Confirmation), CLU (Cross-Lingual Understanding), SSI (Situated Semantic Interpretation), WSI (Wearer-Centered Situational Inference), HII (History-Conditioned Intent Inference), AIP (Action-Trajectory Intent Prediction), STS (Structured Text Synthesis), STN (Spatio-Temporal Navigation Reasoning), GOC (Goal-Oriented Comparison), VTD (Viewpoint Trajectory Description).}
\label{tab:sas_results}
\centering
\resizebox{\textwidth}{!}{
\begin{tabular}{l | cc | ccccc | cccccc | c}\Xhline{1pt}
\textbf{Model} & \multicolumn{2}{c|}{\textbf{L1: Perception}} & \multicolumn{5}{c|}{\textbf{L2: Understanding}} & \multicolumn{6}{c|}{\textbf{L3: Wearer-Grounded Reasoning}} & \textbf{Overall} \\
 & \textbf{L1.1} & \textbf{L1.2} & \textbf{L2.1} & \textbf{L2.2} & \textbf{L2.3} & \textbf{L2.4} & \textbf{L2.5} & \textbf{L3.1} & \textbf{L3.2} & \textbf{L3.3} & \textbf{L3.4} & \textbf{L3.5} & \textbf{L3.6} & \textbf{AVG} \\
 & \textbf{DTR} & \textbf{DTL} & \textbf{QTE} & \textbf{CF} & \textbf{CLU} & \textbf{SSI} & \textbf{WSI} & \textbf{HII} & \textbf{AIP} & \textbf{STS} & \textbf{STN} & \textbf{GOC} & \textbf{VTD} & \\
\Xhline{1pt}
\multicolumn{15}{l}{\emph{Closed-source Models}} \\
GPT-5.2 & 84.4 & \textbf{74.4} & \underline{69.5} & \textbf{77.3} & \textbf{76.5} & \textbf{76.6} & \textbf{56.1} & \textbf{71.8} & \textbf{70.2} & \textbf{74.1} & \textbf{71.5} & \textbf{71.8} & \textbf{70.4} & \cellcolor{gray!15}{\underline{72.6}} \\
Gemini-3-Pro & \underline{84.1} & 75.2 & \textbf{77.0} & \underline{81.7} & \underline{77.8} & \underline{78.7} & \underline{64.7} & \underline{76.4} & \underline{71.4} & \underline{76.9} & \underline{76.4} & \underline{76.7} & \underline{72.4} & \cellcolor{gray!15}{\textbf{76.1}} \\
Claude-Opus-4.5 & 86.8 & \underline{73.1} & 76.1 & 81.5 & 72.5 & 73.4 & 63.6 & 74.0 & 72.8 & 75.3 & 76.6 & 74.4 & 74.3 & \cellcolor{gray!15}{75.0} \\
Qwen3-VL-PLUS & \textbf{89.1} & 72.1 & 73.1 & 77.0 & 74.0 & 72.2 & 65.8 & 71.0 & 69.7 & 74.3 & 73.5 & 73.2 & 70.4 & \cellcolor{gray!15}{73.5} \\
Seed-1.6 & 87.9 & 70.0 & 73.1 & 76.5 & 67.7 & 72.3 & 65.9 & 67.3 & 63.9 & 70.2 & 67.3 & 67.0 & 66.3 & \cellcolor{gray!15}{70.4} \\
\hline
\multicolumn{15}{l}{\emph{Open-source Models}} \\
InternVL3.5-8B & 66.6 & 59.6 & 54.4 & 62.9 & 54.9 & 54.5 & 26.5 & 60.7 & 50.1 & 61.6 & 60.9 & 58.7 & 55.3 & \cellcolor{gray!15}{55.9} \\
InternVL3.5-38B & 79.2 & 68.7 & 54.9 & 55.5 & 58.4 & 54.6 & 36.9 & 51.2 & 51.9 & 54.6 & 51.1 & 52.8 & 59.3 & \cellcolor{gray!15}{56.1} \\
Qwen2.5-VL-3B & 75.0 & 67.7 & 64.8 & 73.6 & 62.0 & 59.7 & 34.7 & 62.0 & 64.5 & 61.2 & 62.1 & 61.0 & 54.9 & \cellcolor{gray!15}{61.8} \\
Qwen2.5-VL-7B & 79.8 & 65.0 & 65.4 & 74.7 & 69.5 & 65.3 & 30.0 & 68.0 & 69.4 & 68.6 & 68.3 & 65.7 & 82.4 & \cellcolor{gray!15}{67.1} \\
Qwen2.5-VL-32B & 83.2 & 69.2 & 64.1 & 76.7 & 70.0 & 65.3 & 51.7 & 64.2 & 66.5 & 67.2 & 65.9 & 65.2 & 64.5 & \cellcolor{gray!15}{67.2} \\
Qwen3-VL-4B & 80.4 & 72.1 & 65.1 & 73.2 & 71.3 & 65.6 & 48.6 & 67.6 & 66.0 & 68.3 & 64.3 & 64.5 & 58.3 & \cellcolor{gray!15}{66.6} \\
Qwen3-VL-8B & 82.4 & 67.7 & 65.1 & 60.8 & 70.6 & 64.5 & 56.9 & 56.1 & 62.8 & 57.1 & 55.6 & 53.1 & 64.7 & \cellcolor{gray!15}{62.9} \\
Qwen3-VL-32B & 83.1 & 70.4 & 66.4 & 76.3 & 71.7 & 69.3 & 57.1 & 65.3 & 68.7 & 69.1 & 69.8 & 66.0 & 69.2 & \cellcolor{gray!15}{69.4} \\
GLM-4.6V-Flash-9B & 80.3 & 63.4 & 60.3 & 44.0 & 57.1 & 57.8 & 42.7 & 46.1 & 51.5 & 47.1 & 47.0 & 46.5 & 44.7 & \cellcolor{gray!15}{53.0} \\
GLM-4.6V-106B-A12B & 83.8 & 73.2 & 64.1 & 73.7 & 71.7 & 66.7 & 51.0 & 60.4 & 62.0 & 64.8 & 61.5 & 59.4 & 62.5 & \cellcolor{gray!15}{65.8} \\
MINICPM-V-4.5 & 78.0 & 64.0 & 64.6 & 74.6 & 52.9 & 61.9 & 45.9 & 61.6 & 59.2 & 61.9 & 60.5 & 57.7 & 42.3 & \cellcolor{gray!15}{60.4} \\
\hline
\multicolumn{15}{l}{\emph{Open-source Video Models}} \\
VideoLLaMA3-7B & 58.4 & 59.9 & 61.3 & 67.9 & 61.1 & 55.5 & 35.6 & 60.9 & 61.8 & 65.3 & 64.3 & 62.1 & 42.5 & \cellcolor{gray!15}{58.2} \\
LLaVA-NeXT-Video & 39.6 & 56.4 & 46.5 & 60.7 & 56.0 & 48.1 & 26.5 & 54.6 & 58.3 & 57.3 & 60.3 & 57.6 & 22.7 & \cellcolor{gray!15}{49.6} \\
VideoR1-7B & 56.8 & 61.0 & 54.8 & 64.4 & 61.4 & 63.9 & 44.1 & 55.2 & 56.1 & 54.4 & 55.9 & 52.7 & 67.9 & \cellcolor{gray!15}{57.6} \\
MOLMO-2-8B & 56.7 & 65.2 & 54.8 & 60.2 & 63.3 & 53.9 & 34.4 & 54.7 & 56.6 & 55.3 & 53.9 & 52.0 & 37.5 & \cellcolor{gray!15}{53.7} \\
\hline
\Xhline{1pt}
\end{tabular}
}
\end{table*}


%%%%%%%%%%%%%%%%%%%%%%%%%%%%%%%%%%%%%%%%%%%%%%%%%%%%%%%%%%%%%%%%%%%%%%%%%%%%%%%
\section{Implementation Details}
\label{sec:impl_details}

\subsection{ Data Construction Pipeline}
\label{sec:impl_mahv}
The OCR tool used in this paper is PaddleOCR~\cite{paddleocr}, and the layout-aware transformation, which reorganizes OCR tokens into reading-order reference text $T_{\text{ref}}$ conditioned on detected bounding boxes, is performed by o4-mini~\cite{openai2025o4mini}; this is a separate stage from scoring and does not interact with the judge. The Fact and Consistency Agents return continuous scores in $[0,1]$; QA pairs whose Fact score falls below the acceptance threshold $\tau=0.7$ or whose Consistency Agent emits a conflict report are routed to the Refinement Agent, which is allowed up to 2 revision passes before any still-failing pair is dropped. The retained pairs then proceed to the human-verification stage. 



\subsection{VerEval Judge Setup}
\label{sec:impl_vereval}
\revion{The VerEval judge operates in two modes. In the \emph{evidence-based mode}, it receives the prediction, question,  and verified ground truth, but not the video frames, and evaluates consistency with this verified evidence record. In the \emph{video-aware mode}, the original video is additionally provided, allowing the judge to re-examine the visual evidence directly. In this paper, all main-table results are obtained in the evidence-based mode using a fixed, locally deployed Qwen3-VL-32B judge~\cite{Qwen3-VL}. We adopt this setting for the 20-model leaderboard because it avoids additional input tokens per video and isolates the evaluated model's perception from a second judge-perception pass. The three scorers each make one call per answer. Decoding uses temperature $0$, a 1{,}024-token cap, and a deterministic seed.}

\subsection{VerEval Scoring Prompts}
\label{scoring_prompt}
\revion{We employ three prompts for model evaluation in the VerEval framework. The {Fact Scoring Prompt} generates a structured factual assessment $S_{fact}=(s_{fact},\mathcal{R}_{fact})$ by detecting discrepancies in key information (e.g., numbers, timestamps, and entities) against video content and ground truth; it reports both a diagnostic score $s_{fact}$ and suggested corrections $\mathcal{R}_{fact}$.} The {Logic Scoring Prompt} assesses the structural integrity of the reasoning process and its alignment with spatio-temporal constraints, outputting the Logic Score (\(S_{logic}\)). The {Consistency Scoring Prompt} jointly reviews the fact judgment, logic judgment, their rationales, and the model answer, producing the Consistency Score (\(S_{consist}\)).
\revion{In the final VerEval score, \(S_{consist}\) refers specifically to the output of the consistency-scoring prompt in Figure~\ref{fig:consistency_score_prompt}. It is not a renamed version of the Fact Agent's structured assessment \(S_{fact}\). We treat its scalar component \(s_{fact}\) as an intermediate diagnostic because the consistency scorer already checks factual evidence while also assessing whether the fact and logic judgments are mutually coherent. Empirically, raw fact scoring showed weaker separation from human judgments than the consistency and logic scores. Therefore, \(s_{fact}\) is excluded from the aggregate, and \(S_{overall}\) is computed only from \(S_{consist}\) and \(S_{logic}\).} The full scoring prompt templates are shown in Figure~\ref{fig:scoring_prompts}.


% \subsection{VerEval Scoring Prompts}
% \label{scoring_prompt}
% We employ three prompts for model evaluation in the VerEval framework. The {Fact Scoring Prompt} generates a diagnostic Fact Score (\(S_{fact}\)) by detecting discrepancies in key information (e.g., numbers, timestamps, entities) against video content and ground truth. The {Logic Scoring Prompt} assesses the structural integrity of the reasoning process and its alignment with spatio-temporal constraints, outputting the Logic Score (\(S_{logic}\)). The {Consistency Scoring Prompt} jointly reviews the fact judgment, logic judgment, their rationales, and the model answer, producing the Consistency Score (\(S_{consist}\)).
% In the final VerEval score, \(S_{consist}\) refers specifically to the output of the consistency-scoring prompt in Figure~\ref{fig:consistency_score_prompt}. It is not a renamed version of \(S_{fact}\). We treat \(S_{fact}\) as an intermediate diagnostic because the consistency scorer already checks factual evidence while also assessing whether the fact and logic judgments are mutually coherent. Empirically, raw fact scoring showed weaker separation from human judgments than the consistency and logic scores. Therefore, \(S_{fact}\) is excluded from the aggregate, and \(S_{overall}\) is computed only from \(S_{consist}\) and \(S_{logic}\). The full scoring prompt templates are shown in Figure~\ref{fig:scoring_prompts}.

\begin{figure*}[p]
  \centering
  \begin{subfigure}{0.76\textwidth}
    \centering
    \includegraphics[width=\linewidth]{figures/fact_score_prompt.pdf}
    \caption{Fact scoring prompt.}
    \label{fig:fact_score_prompt}
  \end{subfigure}
  \vspace{0.4em}

  \begin{subfigure}{0.76\textwidth}
    \centering
    \includegraphics[width=\linewidth]{figures/logic_score_prompt.pdf}
    \caption{Logic scoring prompt.}
    \label{fig:logic_score_prompt}
  \end{subfigure}
  \vspace{0.4em}

  \begin{subfigure}{0.76\textwidth}
    \centering
    \includegraphics[width=\linewidth]{figures/consistency_score_prompt.pdf}
    \caption{Consistency scoring prompt.}
    \label{fig:consistency_score_prompt}
  \end{subfigure}
  \caption{Scoring prompt templates used in VerEval.}
  \label{fig:scoring_prompts}
\end{figure*}



%%%%%%%%%%%%%%%%%%%%%%%%%%%%%%%%%%%%%%%%%%%%%%%%%%%%%%%%%%%%%%%%%%%%%%%%%%%%%%%
\section{Additional Analyses}
\label{sec:additional}

\subsection{MAH-V Quality Control Ablation}
\label{sec:ablation_mahv}
We perform an ablation analysis over the three quality control agents to validate their effectiveness, as shown in Table~\ref{ablation}. For each ablation, we sampled 1{,}391 QA pairs and calculated the average human acceptance over three trained annotators. Removing any quality control module decreases human acceptance: removing the Fact Checker causes the largest drop ($-9.0$ percentage points, to 86.5\%), followed by the Consistency Agent ($-7.1$ points) and the Logic Checker ($-2.8$ points). Removing all agents further degrades acceptance to 84.1\%.

\begin{table}[t!]
\centering
\caption{Ablation study evaluating the impact of each quality control agent on the acceptability of the synthesized tasks.}
\label{ablation}
\resizebox{\linewidth}{!}{
\begin{tabular}{ccc|c}
\toprule
\textbf{Fact Checker} & \textbf{Logical Checker} & \textbf{Consistency Agent} & \makecell{\textbf{Human} \\ \textbf{Acceptance}} \\
\midrule
\cross & \cross & \cross & 84.1\% \\
\cross & \tick  & \tick  & 86.5\% \\
\tick  & \cross & \tick  & 92.7\% \\
\tick  & \tick  & \cross & 88.4\% \\
\tick  & \tick  & \tick  & 95.5\% \\
\bottomrule
\end{tabular}
}
\vspace{-3mm}
\end{table}

% Complementing the end-to-end MAH-V SFT comparison in main-paper Table~\ref{tab:qwen_ablation}, we additionally ablate the layout-aware OCR transformation in Table~\ref{tab:ocr_layout_ablation}. Although the overall gap is modest, removing layout-aware reordering lowers the overall score by 1.7 points and causes a large drop on L1.2 Dynamic Text Localization (19.6~$\to$~7.3), confirming that spatial structure is especially important for text grounding.

\subsection{\revion{Layout-Aware OCR Ablation}}
\label{sec:ocr_layout_ablation}

\revion{To isolate the contribution of spatial layout information, we compare layout-aware OCR with OCR tokens without layout-based reordering in Table~\ref{tab:ocr_layout_ablation}. Removing layout-aware reordering reduces the overall score by 1.7 points, demonstrating that preserving spatial structure is particularly important.}



\begin{table}[t!]
\centering
\caption{OCR layout ablation. Scores are reported by task level and overall average.}
\label{tab:ocr_layout_ablation}
\small
\begin{tabular}{lcccc}
\toprule
\textbf{Input} & \textbf{L1} & \textbf{L2} & \textbf{L3} & \textbf{Overall} \\
\midrule
OCR+Layout & \textbf{33.8} & \textbf{37.6} & 23.4 & \textbf{30.4} \\
OCR-NoLayout & 27.2 & 35.6 & \textbf{23.6} & 28.7 \\
\bottomrule
\end{tabular}
\end{table}

\subsection{VerEval Weighting Ablation}
\label{sec:vereval_weighting_ablation}
We select the final VerEval aggregation weight by maximizing alignment with human judgments on held-out annotated samples. Since \(S_{consist}\) is produced by the consistency prompt and subsumes factual conflicts as well as fact--logic coherence, it serves as the primary term; \(S_{logic}\) provides a secondary penalty for flawed reasoning chains. We report AUC (area under the ROC curve), which measures how well a metric ranks human-labeled \textit{Correct} answers above \textit{Wrong} answers; higher values indicate stronger agreement with human judgments. Table~\ref{tab:vereval_weighting_ablation} shows that the \(0.9/0.1\) consistency--logic weighting achieves the highest average AUC across annotator sets. The variance across configurations is below 0.002, indicating that VerEval is robust to the exact weighting choice rather than sensitive to parameter tuning.

\begin{table}[t!]
\centering
\caption{VerEval weighting ablation measured by AUC against human annotations.}
\label{tab:vereval_weighting_ablation}
\small
\begin{tabular}{lcccc}
\toprule
\textbf{Weight} & \textbf{Ann-1} & \textbf{Ann-2} & \textbf{Ann-3} & \textbf{Avg.} \\
\midrule
0.5 / 0.5 & 0.827 & 0.827 & 0.837 & 0.830 \\
0.7 / 0.3 & 0.827 & 0.826 & 0.840 & 0.831 \\
0.8 / 0.2 & \textbf{0.828} & \textbf{0.827} & \textbf{0.840} & 0.831 \\
0.9 / 0.1 & \textbf{0.828} & \textbf{0.827} & \textbf{0.840} & \textbf{0.832} \\
\bottomrule
\end{tabular}
\end{table}

\subsection{Model-Size Scaling}
\label{sec:scaling}
Figure~\ref{fig:model_size_scaling} shows scaling effects for Qwen2.5-VL and Qwen3-VL. Qwen2.5-VL shows steady gains across all levels, with L1 rising by 9.9 points and L3 by 17.1 points. Conversely, Qwen3-VL displays non-monotonic patterns: while L1 and L2 scale upward, L3 performance regresses at 8B (42.7~$\to$~39.4) before recovering at 32B. These results indicate that perception scales predictably, whereas understanding and reasoning exhibit instability at scales.

\subsection{Video-Pretrained Models}
\label{sec:video_native}
Video-native models show mixed results, with most underperforming image-based counterparts. LLaVA-NeXT-Video exhibits catastrophic failure (15.1 overall), while VideoLLaMA3-7B (40.9) and MOLMO-2-8B (35.6) trail behind similarly-sized image models like Qwen3-VL-4B (51.9). VideoR1-7B achieves competitive performance (50.4) with notable strength in L3.2 (62.4), but still struggles with Structured Text Synthesis (L3.3: 24.2) and Viewpoint Trajectory Description (L3.6: 24.6). These results suggest that conventional video pretraining on third-person, motion-centric datasets does not directly transfer to first-person, text-dense wearable scenarios where frame-level text clarity, wearer-view alignment, and semantic consistency matter more than generic motion modeling.


\subsection{\texorpdfstring{\revion{Language-Stratified Analysis}}{Language-Stratified Analysis}}
\label{sec:language_analysis}
\revion{Table~\ref{tab:detailed_data_stats} provides the complete dataset breakdown by QA and scene-text language. As an additional language-stratified model check reported during rebuttal, Qwen2.5-VL-7B obtains similar VerEval scores on Chinese-only and Chinese--English mixed questions (54.5 vs. 53.9; scores are multiplied by 100). This two-group result suggests no large gap between these strata for this model.}

\begin{table}[t]
\centering
\caption{\revion{Available language-stratified VerEval check for Qwen2.5-VL-7B.}}
\label{tab:language_performance}
\small
\begin{tabular}{lc}
\toprule
\textbf{Question language} & \textbf{VerEval} \\
\midrule
Chinese-only & 54.5 \\
Chinese--English mixed & 53.9 \\
\bottomrule
\end{tabular}
\end{table}

%%%%%%%%%%%%%%%%%%%%%%%%%%%%%%%%%%%%%%%%%%%%%%%%%%%%%%%%%%%%%%%%%%%%%%%%%%%%%%%
\section{\texorpdfstring{\revion{Data Collection and Ethics}}{Data Collection and Ethics}}
\label{sec:collection_protocol}
\label{sec:privacy}

\paragraph{\revion{Participants and Consent.}}
\revion{Eight volunteers, all members of the research team, recorded the first-person videos. There was no external recruitment, paid crowdsourcing, or participant payment. Before recording, every participant received the collection and privacy protocol and provided informed consent with full information about the research purpose, data processing, and planned distribution. All collection procedures followed institutional ethical and privacy guidelines.}

\paragraph{\revion{Intent-Driven Collection Protocol.}}
\revion{Recording was restricted to permitted public settings in shopping, transportation, and tourism scenarios. Participants received instructions covering safety, privacy, and desired coverage of text-rich activities, but were not given scripted actions or questions. Instead, they selected genuine activities they intended to perform, such as reading a sign, confirming a location, comparing displayed options, or navigating with public text. Candidate QAs were grounded in each recorder's video evidence, after which the original recorder checked their intent alignment.}

\paragraph{\revion{Expert Validation.}}
\revion{Following MAH-V, three domain experts independently validated each candidate QA pair against the original video, verified reference text, and task definition. Annotators could correct or reject a problematic question or answer rather than merely assign an acceptability label. Final retention required unanimous approval or a 2/3 majority. This correction-and-majority procedure provides a correctness-oriented check, while the separately reported inter-annotator agreement of $\kappa=0.89$ measures the reliability of the validation decisions.}



\paragraph{Face Anonymization.}
\revion{All third-party faces in videos and global images were automatically detected and blurred while preserving the non-face scene context needed for the tasks. Human reviewers subsequently performed frame-level inspection of the processed data. Any incomplete anonymization was reprocessed and checked again before release.}

\paragraph{Data Usage and Distribution.}
{The anonymized dataset is intended exclusively for academic, non-commercial research. Its usage agreement prohibits attempts to reverse anonymization or identify individuals and requires appropriate attribution and responsible handling. Rejected or insufficiently anonymized material is excluded from the release.}



% %%%%%%%%%%%%%%%%%%%%%%%%%%%%%%%%%%%%%%%%%%%%%%%%%%%%%%%%%%%%%%%%%%%%%%%%%%%%%%%
% \section{Reproducibility}
% \label{sec:repro}
% We release all scoring prompt templates used by VerEval in Section~\ref{scoring_prompt}.
% Our codebase, including model inference and evaluation scripts, will be released through an anonymous repository at the camera-ready stage.
% We also include environment specifications and configuration files for each experiment to ensure exact repeatability.

\end{document}
